\documentclass[11pt,a4paper]{article}
\usepackage[T1]{fontenc}
\usepackage{lmodern}
\usepackage[margin=25mm]{geometry}
\usepackage{amsmath,amssymb,amsthm,mathtools,mathrsfs}
\usepackage{booktabs,longtable,array,enumitem,xcolor,microtype,xurl}
\usepackage[colorlinks=true,linkcolor=blue!45!black,citecolor=blue!45!black,urlcolor=blue!50!black]{hyperref}
\usepackage{bookmark,fancyhdr}
\pagestyle{fancy}
\fancyhf{}
\fancyhead[L]{\small Bernoulli, Tornheim, and digamma identities}
\fancyhead[R]{\small ProveIt continuation}
\fancyfoot[C]{\thepage}
\setlength{\headheight}{14pt}
\setlength{\emergencystretch}{3em}
\setlength{\parskip}{3pt}
\setcounter{tocdepth}{2}
\numberwithin{equation}{section}
\theoremstyle{plain}
\newtheorem{theorem}{Theorem}[section]
\newtheorem{lemma}[theorem]{Lemma}
\newtheorem{proposition}[theorem]{Proposition}
\newtheorem{corollary}[theorem]{Corollary}
\newtheorem{conjecture}[theorem]{Conjecture}
\theoremstyle{definition}
\newtheorem{definition}[theorem]{Definition}
\theoremstyle{remark}
\newtheorem{remark}[theorem]{Remark}
\newtheorem{example}[theorem]{Example}
\DeclareMathOperator{\Li}{Li}
\DeclareMathOperator{\Res}{Res}
\DeclareMathOperator{\CT}{CT}
\DeclareMathOperator{\FP}{FP}
\DeclareMathOperator{\NF}{NF}
\DeclareMathOperator{\wt}{wt}
\DeclareMathOperator{\rank}{rank}
\DeclareMathOperator{\Span}{span}
\newcommand{\Q}{\mathbb Q}
\newcommand{\C}{\mathbb C}
\newcommand{\R}{\mathbb R}
\newcommand{\N}{\mathbb N}
\newcommand{\Z}{\mathbb Z}
\newcommand{\dd}{\,\mathrm d}
\newcommand{\src}[1]{\nolinkurl{#1}}
\newcommand{\CM}{\mathcal C}
\newcommand{\eps}{\varepsilon}
\newcommand{\re}{\operatorname{Re}}
\newcommand{\mi}{\mathrm i}
\hypersetup{pdftitle={Bernoulli Determinants, Cubic Tornheim Jets, and Canonical Digamma Integrals},pdfauthor={ProveIt research continuation}}
\title{\textbf{Bernoulli Determinants, Cubic Tornheim Jets,\\and Canonical Digamma Integrals}\\[8pt]
\large Exact harmonic-tail reconstruction, polylogarithmic boundary identities,\\
and finite-part reduction through the seventh power}
\author{Research continuation for the ProveIt programme}
\date{11 October 2026 (UTC)\\[4pt]
\small Repository snapshot \texttt{f7e35be357654198c1546fd40f363694c5bcc2cc}}

\begin{document}
\maketitle
\begin{abstract}
We develop three exact continuations of the current ProveIt
polylogarithm--Hurwitz--Stieltjes programme. First, the conjectured
optimal consecutive moment prefix is proved for every unequal even pair
of harmonic-tail orders with smaller order $2$, $4$, or $6$. For $(2,2b)$,
$b\ge2$, the decisive coefficient determinant is
$(-1)^{b-1}(2-B_{2b-2})/(2^b(2b-2)!)$. A rank-two determinant with an arithmetic nonvanishing proof settles the
order-four family. A modulo-four Bernoulli congruence and an exact
polynomial certificate settle order six. A fixed-size determinant reduction is given
for every even pair. For order two, a row of the inverse
produces an all-order family of absolutely convergent harmonic and
polygamma sums with rational values, together with exact depth-two
polylogarithmic boundary identities. Second, every admissible third directional
derivative of the Tornheim function at the spectral origin is expressed
using two specified higher-depth constants and ordinary zeta derivatives.
Two elementary homogeneous conditions characterize formal cancellation
of these constants; a three-ray example evaluates exactly. A local
holomorphic-kernel theorem extends this structural analysis to every
Taylor degree. Third, we give a terminating integration-by-parts normal
form for differential polynomials, indexed by partitions whose largest
part is repeated. A Gamma-adjoint identity imposes a further parity
restriction. With every endpoint constant retained, the canonical sixth
digamma moment reduces to the height-one Laurent coefficient
$\lambda_5$, one mixed finite-part integral, and ordinary zeta data.
A seventh-degree identity and a general coincident-polygamma product
formula complete the explicit applications. All uniform statements have proofs. Exact finite certificates and
independent high-precision diagnostics are supplied separately. The
remaining unequal-order moment-prefix conjecture for $8\le p<q$, arithmetic elimination
of the remaining mixed constants, and the manuscript's $S_6$ and revised
$S_8$ identities remain open.
\end{abstract}
\clearpage
\begingroup
\small
\setlength{\parskip}{0pt}
\tableofcontents
\endgroup
\clearpage

% BEGIN sections/01_scope.tex
\section{Scope, provenance, and the exact questions continued}
\label{sec:scope}

The starting point is the current collective manuscript
\emph{Polylogarithms and their Arithmetic Bridges} and the six
archives in the repository's incoming directory at the title-page
commit \cite{Repo}. The source was obtained as a pinned Git checkout;
the incoming archives were extracted into separate directories outside
that checkout. The research package is an addition for later
integration. It does not change the canonical manuscript or the
incoming archive contents.

The two newest digamma reports are particularly relevant.
\emph{Canonical Digamma and Fifth Moment} proves a finite, but
generally longer, initial-moment reconstruction theorem and conjectures
that the first $q$ moments already suffice for unequal even orders
\cite[Sections 7 and 9, \texttt{conj:firstq}]{CanonicalFifth}.
\emph{Higher Digamma Powers and Ordered Gamma Sums} asks for an
organized sixth-power reduction and for a differential-polynomial
integration calculus \cite[Further exact research questions]{HigherPowers}.
The previous coincident/directional-zeta report supplies the complete
quadratic ray formula for Tornheim functions, while leaving cubic
arithmetic data unresolved \cite[Section 6]{DirectionalPrior}.

\subsection{Main contributions and their limits}

The first contribution resolves an infinite subfamily of a stated
conjecture, rather than adding another finite rank experiment.
For the pair $(2,2b)$, an explicit Pad\'e calculation evaluates the
coefficient determinant and a rational extraction row. Classical
Bernoulli Hankel machinery is reproved as needed; the mixed matrix
perturbation and its application to the moment conjecture are the
new continuation \cite{DilcherJiu2020}. The resulting sums and the
polylogarithmic boundary family are identities at every $b\ge2$.
A rank-two version proves the entire $(4,2b)$ family. A congruence
modulo four and a finite polynomial certificate prove $(6,2b)$ as well. For every
even pair $(2a,2b)$ the remaining determinant is reduced explicitly to
size $a$; the cases $8\le p<q$ remain unresolved here.

The second contribution is an exact cancellation theorem for third
spectral derivatives. The spectral origin is not a regular point of
the three-variable Tornheim function. We first isolate its polar
part, then differentiate specified holomorphic ray restrictions.
The elementary axis and quasi-shuffle restrictions determine enough
of the remaining Taylor coefficients to yield a complete two-coordinate
formula and explicit combinations containing ordinary zeta derivatives
alone. The formal ambiguity theorem states precisely what these
particular local restrictions leave undetermined. It is not a
transcendence or arithmetic-independence theorem.

The third contribution answers the algebraic part of the proposed
differential-polynomial classification. Its partition basis gives
an exact finite algorithm for reduction modulo total derivatives.
This is a constructive specialization of classical differential
algebra \cite{OlverShakiban1978}, with a new application to the current
Gamma-normalized finite-part problem. The normalizations at degrees
six and seven are computed explicitly. A remaining mixed integral is displayed,
not silently declared independent or replaced by an unevaluated
quantity under a new name.

The principal status changes are summarized in
Table~\ref{tab:status}. All references to novelty in this article
mean an addition relative to the inspected ProveIt material.
A complete historical priority claim would require a broader
comparison with all specializations of the classical literature.

\begin{table}[htbp]
\centering\small
\caption{Source questions and the result established here.}
\label{tab:status}
\begin{tabular}{@{}p{.29\textwidth}p{.42\textwidth}p{.22\textwidth}@{}}
\toprule
Question & Proved result & Remaining scope\\
\midrule
Optimal first-$q$ prefix & Complete proofs for smaller order $2$, $4$, or $6$; a fixed-size determinant for every even pair & Unequal $8\le p<q$\\[5pt]
Higher Tornheim ray derivatives & All cubic directions; exact cancellation criterion and an all-degree local kernel & Arithmetic values of the two cubic coordinates\\[5pt]
Differential digamma classification & Canonical partition basis, bidegree dimensions, Gamma parity and endpoints & The image after additional spectral relations\\[5pt]
Sixth and seventh moments & Canonical sixth formula with one mixed moment; explicit seventh formula & Spectral evaluation of the mixed integrals\\[5pt]
Short Gaussian evaluations & No status change & $S_6$ and revised $S_8$\\
\bottomrule
\end{tabular}
\end{table}

\subsection{Conventions that affect the constants}

For $r\ge1$, set $H_n^{(r)}=\sum_{k=1}^n k^{-r}$ and
$H_0^{(r)}=0$. Bernoulli numbers and polynomials use
\[
 \frac{te^{xt}}{e^t-1}=\sum_{n\ge0}B_n(x)\frac{t^n}{n!},
 \qquad B_n=B_n(0),\qquad B_1=-\frac12.
\]
The digamma function is $\psi(x)=\Gamma'(x)/\Gamma(x)$.
The generalized Stieltjes convention is
\begin{equation}\label{eq:stieltjes-convention}
 \zeta(1+u,a)=\frac1u+
       \sum_{j\ge0}\frac{(-1)^j\gamma_j(a)}{j!}u^j,
 \qquad \re a>0.
\end{equation}
We abbreviate $\gamma=\gamma_0(1)$ and
$\gamma_j=\gamma_j(1)$. Derivatives $\zeta^{(j)}(s)$ are always
with respect to the spectral variable, not to a Hurwitz shift.
Standard Gamma, polygamma, and Hurwitz formulas are used with these
normalizations \cite{DLMF}.

If $f$ is meromorphic near $0$ and regular on $(0,1]$, then
\begin{equation}\label{eq:fp-convention}
 \FP\int_0^1 f(x)\dd x
 :=\CT_{\eps\downarrow0}\int_{\eps}^1 f(x)\dd x.
\end{equation}
The constant term discards negative powers of $\eps$ and logarithms
but retains the constant in the original coordinate $x$.
In particular, for a meromorphic primitive $g$,
\begin{equation}\label{eq:endpoint-law-intro}
 \FP\int_0^1 g'(x)\dd x=g(1)-\CT_{x=0}g(x).
\end{equation}
This convention is distinct from the centered asymptotic moment
continuation in Section~\ref{sec:tail-foundation}; the latter is
defined explicitly before it is used. Both conventions eventually
lead to ordinary convergent identities after their stated subtractions.

% END sections/01_scope.tex


% BEGIN sections/02_tail_foundation.tex
\section{Centered harmonic moments and the exact Bernoulli transform}
\label{sec:tail-foundation}

We restate the precise inherited moment formulas needed for the new
determinant. This also makes the continuation independent of a
choice of numerical regularization. Fix an even $q=2b\ge4$ and put
\[
 A_s(n)=\zeta(s)-H_n^{(s)}=\zeta(s,n+1),\qquad x_n=n+\frac12,
\]
\begin{equation}\label{eq:tail-Dirichlet}
 D_{2,q}(s)=\sum_{n\ge0} A_2(n)A_q(n)x_n^{-s},
 \qquad \re s>1-q.
\end{equation}
The displayed half-plane is the domain of absolute convergence.
For later $s$, the notation means its unique meromorphic continuation.
We write $D_r=D_{2,q}(-2r)$ for $r\ge0$.

\subsection{Regularity at every required moment}

The centered Hurwitz expansion, obtained from Euler--Maclaurin, is
\begin{equation}\label{eq:centered-hurwitz-foundation}
 A_s(n)\sim x_n^{1-s}\left\{
 \frac1{s-1}+\sum_{h\ge1}
    \frac{B_{2h}(1/2)(s)_{2h-1}}{(2h)!}x_n^{-2h}\right\}.
\end{equation}
Here $s$ in $A_s$ is a fixed positive even integer, and
$(s)_k$ is the rising factorial. The vanishing of all odd
$B_{2h+1}(1/2)$ explains the parity of this expansion.
Consequently $A_2(n)A_q(n)$ has only powers $x_n^{-q-2h}$.
Subtracting finitely many of them from \eqref{eq:tail-Dirichlet}
continues the function to any prescribed left half-plane. Its only
possible poles occur at $s=1-q-2h$, which are negative odd integers.
Every $D_r$ is therefore regular.

This argument also proves a useful operational rule. If $P$ is an
even polynomial and $R$ is the nonnegative-power part of the
asymptotic series for $P(x_n)A_2(n)A_q(n)$, then
\begin{equation}\label{eq:ordinary-fp-foundation}
 \sum_{n\ge0}\{P(x_n)A_2(n)A_q(n)-R(x_n)\}
 =\sum_r [X^{2r}]P(X)\,D_r.
\end{equation}
The left series is absolutely convergent. Each subtracted monomial
has continued sum $\zeta(-2h,1/2)=0$; this follows either from
$\zeta(s,1/2)=(2^s-1)\zeta(s)$ or from the odd Bernoulli polynomial
at $1/2$. Thus the analytic continuation and the explicitly subtracted
ordinary sum have exactly the same constant.

\subsection{The fixed formal coordinate system}

Write decreasing-index multiple zeta values as
\[
 \zeta(a,c)=\sum_{n>m\ge1}\frac1{n^a m^c},\qquad a\ge2,\ c\ge1.
\]
Define the initial blocks by
\begin{align}
 C_0={}&\zeta(2,q-1)+\zeta(q,1)+\zeta(q+1),
 \label{eq:C0-foundation}\\
 C_j={}&\zeta(2,q-2j-1)+\frac12\zeta(q+1-2j)
   +\frac1{2j}\sum_{h=1}^{j}\binom{2j}{2h}
                  B_{2j-2h}\zeta(q-2h),
 \quad 1\le j<b.\label{eq:Cinitial-foundation}
\end{align}
For $j\ge b$, define $C_j=\sum_{k=0}^{b-1}a_{jk}Z_k$, where
\begin{equation}\label{eq:coordinates-foundation}
 Z_0=\zeta(0)=-\frac12,\qquad Z_k=\zeta(2k)\quad(1\le k<b),
\end{equation}
and, with $d=2j+1$, $e=d-q+2k-1$,
\begin{equation}\label{eq:entries-foundation}
 a_{jk}=B_e\left\{
 \frac{\binom{d-q+1}{2-2k}}{d-q+1}
 +\frac{\binom{d-1}{q-2k}}{d-1}\right\}.
\end{equation}
A binomial summand in this formula is omitted unless its lower
index is at least $2$ and no greater than its upper index.
In particular the first summand occurs only when $k=0$.
This is part of the formula, not a convention that may be replaced
by $\binom u0=1$.

\begin{proposition}[Inherited two-tail transform]
\label{prop:tail-transform}
For every $r,j\ge0$,
\begin{align}
 D_r&=\sum_{j=0}^r
       \frac{\binom{2r}{2j}}{2j+1}B_{2r-2j}(1/2)C_j,
       \label{eq:forward-foundation}\\
 C_j&=\sum_{r=0}^j
       \binom{2j+1}{2r}2^{-2(j-r)}D_r.
       \label{eq:inverse-foundation}
\end{align}
Hence all moments lie in the rational span of the $q$ specified
formal labels
\[
 C_0,\ldots,C_{b-1},\qquad Z_0,\ldots,Z_{b-1}.
\]
\end{proposition}

These are the $p=2$ cases of the already proved two-tail formulas
reproduced in \cite[Section 7]{CanonicalFifth}. We give their finite
mechanism to fix the cutoffs and constants. Put
\[
 S_r(N)=\sum_{n=0}^{N-1}(n+1/2)^{2r}
       =\sum_{j=0}^r\frac{\binom{2r}{2j}}{2j+1}
                        B_{2r-2j}(1/2)N^{2j+1}.
\]
This is an odd polynomial with zero constant term. The exact identity
$A_s(k-1)=A_s(k)+k^{-s}$ and finite summation by parts express $D_r$
as the same linear combination of constant terms of
\begin{align}\label{eq:finite-E-foundation}
 E_d(N)={}&N^dA_2(N)A_q(N)
 +\sum_{k=1}^N k^{d-q}A_2(k)
 +\sum_{k=1}^N k^{d-2}A_q(k)
 +\sum_{k=1}^N k^{d-q-2}.
\end{align}
The sign of the first term and the inclusive tail cutoff both
matter. When $d<q$, the first ordered sum tends to
$\zeta(2,q-d)$. For a positive integer exponent $m$, the exact reduction
\begin{equation}\label{eq:finite-Faulhaber-foundation}
 \sum_{k=1}^N k^m A_s(k)
 =A_s(N)F_m(N)+\sum_{\ell=1}^N\frac{F_m(\ell-1)}{\ell^s},
 \quad F_m(N)=\frac{B_{m+1}(N+1)-B_{m+1}}{m+1},
\end{equation}
reduces all remaining sums to finite Bernoulli polynomials and
ordinary power sums. Their constant terms are
\eqref{eq:C0-foundation}--\eqref{eq:entries-foundation}; at the
transition $d=q+1$ the two $-\gamma/2$ contributions cancel the
$+\gamma$ from the final harmonic sum in
\eqref{eq:finite-E-foundation}. The remaining boundary products give
the $Z_0$ column. This is the inherited finite coefficient calculation,
not an assumption of arithmetic independence. Finally the standard
Bernoulli generating identity inverts the odd-polynomial transform,
giving \eqref{eq:inverse-foundation}.

The existing report proves that a longer specified moment prefix
has full formal rank. The next section determines the exact first-$q$
minor for every $p=2$, without any numerical nonvanishing assumption.

% END sections/02_tail_foundation.tex


% BEGIN sections/03_moment_rank.tex
% Research fragment for integration into the parent article.
% All assertions in this fragment have a self-contained algebraic proof.
% The notation C_j,D_r,a_{jk} is the notation of the source tail-rank section.

\section{The optimal moment prefix for every pair \texorpdfstring{$(2,q)$}{(2,q)}}
\label{sec:rank-two}

The source conjecture asks whether the first $q$ centered moments have
full formal rank for every pair of even orders $2\leq p<q$.
We prove the whole infinite subfamily $p=2$, with an exact determinant,
an explicit rational extraction row, and an ordinary convergent
polygamma-tail identity.  The order-four and order-six families are proved below; unequal pairs $8\leq p<q$
remain open here.

The Pad\'e and Hankel ingredients used below are classical.  In
particular, the normalized shifted Bernoulli Hankel determinant is
recorded by Dilcher and Jiu~\cite[\S7.3, Eq.~(7.12)]{DilcherJiu2020},
who credit the evaluation to Krattenthaler.  We reprove the needed
formula to make the normalization and the inverse row explicit.
The new step relative to the source tail-moment report is the
rank-one perturbation calculation giving the factor $2-B_{q-2}$,
and its consequences for the optimal prefix and rational tail sums.

Write $q=2b$, where $b\geq2$, and put
\[
 M_b=(a_{b+t,k})_{0\leq t,k<b}.
\]
The entry convention is exactly the source convention: a binomial
summand is omitted if its lower index is less than $2$.  Thus, when
$p=2$, the first binomial summand occurs only in column $k=0$.

\begin{theorem}[Exact determinant for the entire order-two family]
\label{thm:rank-two}
For every integer $b\geq2$,
\begin{equation}\label{eq:rank-two-det}
 \det M_b=
 \frac{(-1)^{b-1}\bigl(2-B_{2b-2}\bigr)}
 {2^b(2b-2)!}\ne0.
\end{equation}
Consequently, for $p=2$ and $q=2b$, the first $q$ centered moments
$D_0,\ldots,D_{q-1}$ have full rank $q$ on the source's formal
coordinate list
\[
 C_0,\ldots,C_{b-1},\quad Z_0,\ldots,Z_{b-1},
 \qquad Z_k=\zeta(2k),\quad Z_0=-\tfrac12.
\]
The determinant of their full coefficient matrix, in this coordinate
order, is
\begin{equation}\label{eq:rank-two-full-det}
 \frac{(-1)^{b-1}\bigl(2-B_{2b-2}\bigr)}
 {2^b(2b-2)!\,(4b-1)!!}.
\end{equation}
This is a formal-rank assertion; it does not assert arithmetic
independence of the numerical coordinates.
\end{theorem}

\subsection{A finite Pad\'e calculation, proved explicitly}

We use Bernoulli numbers with generating function
\[
 \frac{t}{e^t-1}=\sum_{r\geq0}B_r\frac{t^r}{r!},
 \qquad B_1=-\frac12,
\]
and write
\[
 \mu_r=\frac{B_{2r}}{(2r)!},\qquad
 F(z)=\sum_{r\geq0}\mu_rz^r
     =\frac{\sqrt z}{2}\coth\frac{\sqrt z}{2}.
\]
The expression involving $\sqrt z$ is an even power series and hence
defines a single-valued holomorphic germ at $z=0$.

\begin{lemma}[Explicit exponential Pad\'e identity]
\label{lem:exp-pade-rank}
For $N\geq0$, define
\[
 P_N(t)=\sum_{k=0}^N
   \frac{(2N-k)!\,N!}{(2N)!\,k!\,(N-k)!}\,t^k.
\]
Then
\begin{equation}\label{eq:exp-pade-rank}
 e^tP_N(-t)-P_N(t)
 =\frac{(-1)^N(N!)^2}{(2N)!(2N+1)!}\,t^{2N+1}
  +O(t^{2N+2}).
\end{equation}
\end{lemma}

\begin{proof}
For $0\leq r\leq2N$, the coefficient of $t^r$ in
$e^tP_N(-t)$ equals
\[
 \frac{(2N-r)!}{(2N)!}
 \sum_{k=0}^N(-1)^k\binom Nk
                 \binom{2N-k}{2N-r}.
\]
The sum is the coefficient of $x^{2N-r}$ in
\[
 \sum_{k=0}^N(-1)^k\binom Nk(1+x)^{2N-k}
 =x^N(1+x)^N.
\]
It is therefore $\binom Nr$ for $r\leq N$ and zero for $N<r\leq2N$.
These are precisely the coefficients of $P_N(t)$ through degree
$2N$.  The next coefficient is
\[
 \frac1{(2N)!}\sum_{k=0}^N
       \frac{(-1)^k\binom Nk}{2N+1-k}
 =\frac1{(2N)!}\int_0^1x^N(x-1)^N\,dx,
\]
which gives the stated beta-integral value.
\end{proof}

For $n\geq0$, take $N=2n+1$ and form the polynomials
\begin{align}
 q_n(t^2)&=\frac{P_{2n+1}(t)-P_{2n+1}(-t)}t,
 \\
 r_n(t^2)&=\frac{P_{2n+1}(t)+P_{2n+1}(-t)}2.
\end{align}
Both have degree $n$ and $q_n(0)=1$.  Direct simplification gives
\begin{equation}\label{eq:pade-coth-rank}
 F(z)q_n(z)-r_n(z)
 =h_nz^{2n+1}+O(z^{2n+2}),
 \qquad
 h_n=\frac{((2n+1)!)^2}{(4n+2)!(4n+3)!}.
\end{equation}
Indeed, after setting $z=t^2$, the left side is exactly
\[
 -\frac{e^tP_{2n+1}(-t)-P_{2n+1}(t)}{e^t-1}.
\]
This proves both the vanishing coefficients and the leading
coefficient in \eqref{eq:pade-coth-rank}; no normality theorem for
Pad\'e approximants is being assumed.

Write
\begin{equation}\label{eq:q-explicit-rank}
 q_n(z)=\sum_{\ell=0}^nq_{n,\ell}z^\ell,
 \qquad
 q_{n,\ell}=
 \frac{2(4n+1-2\ell)!\,(2n+1)!}
 {(4n+2)!\,(2\ell+1)!\,(2n-2\ell)!},
\end{equation}
and reverse its coefficients:
\[
 \pi_n(x)=x^nq_n(1/x).
\]
This is a monic polynomial of degree $n$.
Let $\mathcal L$ be the linear functional on polynomials specified by
$\mathcal L(x^r)=\mu_r$.  Comparing coefficients in
\eqref{eq:pade-coth-rank} proves
\begin{align}
 \mathcal L(x^k\pi_n(x))&=0 &&(1\leq k\leq n),
       \label{eq:orthogonality-rank}\\
 \mathcal L(\pi_n)&=c_n,
 &c_n&=\frac{(2n+1)(2n+2)!}{(4n+2)!},
       \label{eq:constant-moment-rank}\\
 \mathcal L(x^{n+1}\pi_n)&=h_n.
       \label{eq:shifted-norm-rank}
\end{align}
Here \eqref{eq:constant-moment-rank} is the coefficient of $z^n$
in $r_n(z)$.

\subsection{The normalized Bernoulli Hankel determinant}

Put
\[
 H_b=(\mu_{i+j})_{0\leq i,j<b},\qquad
 \Delta_n=\det(\mu_{i+j+1})_{0\leq i,j<n},\qquad\Delta_0=1.
\]
The functional $\mathcal N(g)=\mathcal L(xg)$ makes the monic
polynomials $\pi_0,\pi_1,\ldots$ orthogonal by
\eqref{eq:orthogonality-rank}.  Their squared norms are $h_j$:
the polynomial $\pi_j(x)-x^j$ has degree less than $j$, so
\[
 \mathcal N(\pi_j^2)=\mathcal L(x^{j+1}\pi_j)=h_j.
\]
The triangular change from monomials to these monic polynomials has
determinant $1$.  Therefore
\begin{equation}\label{eq:delta-product-rank}
 \Delta_n=\prod_{j=0}^{n-1}
       \frac{((2j+1)!)^2}{(4j+2)!(4j+3)!}.
\end{equation}
In $H_b$, replace its last row by the linear combination of rows
whose coefficients are those of $\pi_{b-1}$.  This operation leaves
the determinant unchanged because $\pi_{b-1}$ is monic.
The resulting last row is $(c_{b-1},0,\ldots,0)$, whence
\begin{equation}\label{eq:H-product-rank}
 \det H_b=(-1)^{b-1}c_{b-1}\Delta_{b-1}.
\end{equation}
In particular, the Bernoulli Hankel matrix is nonsingular.

There is also a useful analytic explanation of this nonvanishing.
Euler's even-zeta identity gives, for $r\geq1$,
\[
 (-1)^{r-1}\mu_r
 =2\sum_{m\geq1}(2\pi m)^{-2r}.
\]
After conjugation by $\operatorname{diag}(1,-1,1,\ldots)$, the
bottom-right $(b-1)$-square block of $H_b$ is the negative of a
positive definite moment matrix.  Positivity is strict because a
nonzero polynomial cannot vanish at the infinitely many distinct
points $(2\pi m)^{-2}$.  The top-left entry remains $1$, so its
Schur complement is positive.  This yields the sign
$(-1)^{b-1}$ and invertibility independently of the exact product.
The exceptional initial coefficient $\mu_0=1$ is retained; replacing
it by a positive Stieltjes mass would change the problem.

\subsection{The mixed matrix is a rank-one perturbation}

The source coefficient formula, specialized to $p=2$, is
\begin{equation}\label{eq:rank-two-entry}
 (M_b)_{tk}
 =\frac{(2b+2t-1)!}{(2b-2k)!}\,\mu_{t+k}
   +\mathbf1_{k=0}\frac{2t+1}{2}B_{2t}.
\end{equation}
Define diagonal matrices
\[
 R_b=\operatorname{diag}_{0\leq t<b}(2b+2t-1)!,\qquad
 E_b=\operatorname{diag}_{0\leq k<b}\frac1{(2b-2k)!},
\]
and a column vector $u=(u_t)_{0\leq t<b}$ by
\[
 u_t=\mu_t\rho_t,\qquad
 \rho_t=\frac{(2b)!\,(2t+1)!}{2(2b+2t-1)!}.
\]
Then
\begin{equation}\label{eq:rank-one-factor}
 M_b=R_b(H_b+u e_0^{\mathsf T})E_b,
\end{equation}
where $e_0=(1,0,\ldots,0)^{\mathsf T}$.

Set $n=b-1$.  By \eqref{eq:orthogonality-rank} and
\eqref{eq:constant-moment-rank}, the first row of $H_b^{-1}$ is the
coefficient row of $\pi_n/c_n$.  Consequently
\begin{equation}\label{eq:contraction-rank}
 e_0^{\mathsf T}H_b^{-1}u
 =\frac1{c_{b-1}}\sum_{t=0}^{b-1}q_{b-1,b-1-t}\,\mu_t\rho_t.
\end{equation}
The factorials in \eqref{eq:q-explicit-rank} cancel to give
\begin{align}
 \frac1{c_{b-1}}\sum_{t=0}^{b-1}
       q_{b-1,b-1-t}\mu_t\rho_t
 &=\frac1{2b-1}\sum_{t=0}^{b-1}
       \binom{2b-1}{2t}(2t+1)B_{2t}\\
 &=1-B_{2b-2}.
       \label{eq:bernoulli-contraction-rank}
\end{align}
For clarity, the last equality is an elementary Bernoulli identity.
If $N=2b-1\geq3$, then
\[
 \sum_{k=0}^N\binom NkB_k=B_N(1)=0,
 \qquad
 \sum_{k=0}^Nk\binom NkB_k=-NB_{N-1}.
\]
Removing the $k=1$ terms and using $B_k=0$ for odd $k>1$ shows
\[
 \sum_{\substack{0\leq k<N\\ k\text{ even}}}
      (k+1)\binom NkB_k=N(1-B_{N-1}),
\]
as required.  The condition $b\geq2$ matters: this identity is not
the corresponding formula for $N=1$.

The matrix determinant lemma now yields
\[
 \det(H_b+u e_0^{\mathsf T})
 =\det H_b\bigl(2-B_{2b-2}\bigr).
\]
It remains to simplify the factorial factor.  Equations
\eqref{eq:delta-product-rank}--\eqref{eq:H-product-rank} give
\begin{equation}\label{eq:factorial-cancel-rank}
 \det R_b\det E_b\det H_b
 =\frac{(-1)^{b-1}}{2^b(2b-2)!}.
\end{equation}
One quick verification, also a proof for every $b$, is induction:
the left side is $1/2$ at $b=1$ and the quotient of its values at
$b$ and $b-1$ is
\[
 -\frac1{4(b-1)(2b-3)}.
\]
This is exactly the quotient on the right side of
\eqref{eq:factorial-cancel-rank}.

For completeness, $B_{2m}$ cannot equal $2$ for any $m\geq1$ by an
elementary parity argument; the full von Staudt--Clausen theorem is
unnecessary.  Work in the ring $\mathbb Z_{(2)}$ of rational numbers
with odd denominator.  The Bernoulli recurrence gives
\[
 (2m+1)(2B_{2m})
 =2m-1-\sum_{j=1}^{m-1}\binom{2m+1}{2j}(2B_{2j}).
\]
Inductively $2B_{2j}\in\mathbb Z_{(2)}$ and is congruent to $1$
modulo $2$.  The binomial sum in the last display is even modulo
$2$, since
\[
 \sum_{j=1}^{m-1}\binom{2m+1}{2j}
 =2^{2m}-2m-2.
\]
Thus $2B_{2m}$ is also congruent to $1$ modulo $2$.  In particular
its $2$-adic valuation is $0$, whereas $2\cdot2$ is even.
This proves $2-B_{2m}\ne0$ and completes
\eqref{eq:rank-two-det}.  Finally the forward Bernoulli transform
has diagonal entries $1/(2r+1)$, giving
\eqref{eq:rank-two-full-det}.

\subsection{An all-order rational extraction identity}

The same calculation identifies one entire row of the inverse,
without an unspecified matrix inversion.

\begin{theorem}[Explicit extraction of the rational coordinate]
\label{thm:rational-row-two}
For $b\geq2$,
\begin{equation}\label{eq:rational-block-two}
 \sum_{t=0}^{b-1}\binom{2b-1}{2t}\,C_{b+t}
 =-\frac{(2b-1)(2-B_{2b-2})}{4}.
\end{equation}
Equivalently, with
\begin{equation}\label{eq:moment-weights-two}
 w_{b,r}=\sum_{\substack{0\leq t<b\\r\leq b+t}}
 \binom{2b-1}{2t}\binom{2b+2t+1}{2r}
       2^{-2(b+t-r)},\qquad 0\leq r\leq2b-1,
\end{equation}
the centered moments satisfy
\begin{equation}\label{eq:rational-moment-two}
 \sum_{r=0}^{2b-1}w_{b,r}D_r
 =-\frac{(2b-1)(2-B_{2b-2})}{4}.
\end{equation}
\end{theorem}

\begin{proof}
The rank-one inverse formula and
\eqref{eq:bernoulli-contraction-rank} show that
\[
 e_0^{\mathsf T}(H_b+u e_0^{\mathsf T})^{-1}
 =\frac{e_0^{\mathsf T}H_b^{-1}}{2-B_{2b-2}}.
\]
Insert the coefficient formula for $q_{b-1}$ and the diagonal
factors in \eqref{eq:rank-one-factor}.  The resulting first row is
\begin{equation}\label{eq:inverse-first-row-two}
 (M_b^{-1})_{0t}=
 \frac{2}{(2b-1)(2-B_{2b-2})}\binom{2b-1}{2t}.
\end{equation}
Multiplication by $(C_b,\ldots,C_{2b-1})^{\mathsf T}$ recovers
$Z_0=-1/2$, which proves \eqref{eq:rational-block-two}.
Substitution of the source inverse Bernoulli transform
\[
 C_j=\sum_{r=0}^j\binom{2j+1}{2r}2^{-2(j-r)}D_r
\]
gives \eqref{eq:rational-moment-two}.
\end{proof}

As an exact check, for $b=3$ the block identity is
\[
 C_3+10C_4+5C_5=-\frac{61}{24}.
\]
Multiplying its moment form by $1024$ reproduces the source's
$(2,6)$ identity with coefficients
$61,8204,116000,370048,303360,56320$ and right side $-7808/3$.
The source example is therefore a specialization of an all-order
formula, rather than an isolated elimination.

\subsection{An ordinary convergent series at every even order}

The finite-part identity has a direct, absolutely convergent
interpretation with an explicitly specified polynomial subtraction.
Set
\begin{align}
 S_b(y)&=\frac{y^{2b+1}}2
       \bigl((1+y)^{2b-1}+(1-y)^{2b-1}\bigr),
       \label{eq:S-polynomial-two}\\
 W_b(x)&=S_b(x+1/2)-S_b(x-1/2)
        =\sum_{r=0}^{2b-1}w_{b,r}x^{2r}.
       \label{eq:W-polynomial-two}
\end{align}
For an even integer $s\geq2$, define rational coefficients
\[
 c_{s,0}=\frac1{s-1},\qquad
 c_{s,h}=\frac{B_{2h}(1/2)}{(2h)!}(s)_{2h-1}\quad(h\geq1),
\]
where $(s)_r$ is a rising factorial.  Put
\[
 T_{b,h}=\sum_{i=0}^hc_{2,i}c_{2b,h-i},
\]
and define the explicit even polynomial
\begin{equation}\label{eq:subtraction-two}
 R_b(x)=\sum_{r=b}^{2b-1}\ \sum_{h=0}^{r-b}
       w_{b,r}T_{b,h}x^{2(r-b-h)}.
\end{equation}
Only the finitely many coefficients with $h\leq b-1$ occur.

\begin{corollary}[Convergent rational tail sums]
\label{cor:ordinary-two}
Let $x_n=n+1/2$ and $A_s(n)=\zeta(s)-H_n^{(s)}$.  For every
$b\geq2$,
\begin{equation}\label{eq:ordinary-two}
 \sum_{n=0}^{\infty}
 \left\{W_b(x_n)A_2(n)A_{2b}(n)-R_b(x_n)\right\}
 =-\frac{(2b-1)(2-B_{2b-2})}{4}.
\end{equation}
The summand is $O(n^{-2})$, so the series is absolutely convergent.
Equivalently, one may replace the tail product by
\[
 A_2(n)A_{2b}(n)
 =\frac{\psi'(n+1)\psi^{(2b-1)}(n+1)}{(2b-1)!}.
\]
\end{corollary}

\begin{proof}
The centered Hurwitz expansion gives
\[
 A_s(n)\sim x_n^{1-s}\sum_{h\geq0}c_{s,h}x_n^{-2h},
 \qquad
 A_2(n)A_{2b}(n)\sim
       x_n^{-2b}\sum_{h\geq0}T_{b,h}x_n^{-2h}.
\]
The polynomial $R_b$ in \eqref{eq:subtraction-two} is precisely the
sum of all nonnegative powers in the product with $W_b$.
After its subtraction the next possible term is $O(x_n^{-2})$.
Every removed term has continued sum zero because
$\zeta(-2r,1/2)=0$ for $r\geq0$.  Thus the absolutely convergent
sum equals the finite-part moment combination in
\eqref{eq:rational-moment-two}.  Alternatively, the partial sum of
each removed centered even power is an odd polynomial in the upper
integer cutoff with no constant term, which proves the same
constant-term equality without interchanging a divergent series.
The polygamma expression is the standard positive-even-order
identity $A_s(n)=\psi^{(s-1)}(n+1)/(s-1)!$.
\end{proof}

For example, when $b=2$, multiplication by $64$ gives the compact
ordinary identity
\begin{equation}\label{eq:ordinary-24}
 \sum_{n=0}^{\infty}\left\{
 (7+412x_n^2+2000x_n^4+1344x_n^6)A_2(n)A_4(n)
 -448x_n^2-\frac{1216}{3}
 \right\}=-88.
\end{equation}

An independently reconstructed example at $b=4$ is also useful for
checking the size and sign of the subtraction. Define
\begin{align}
 \mathsf P_4(X)={}&547+129756X^2+3657584X^4+27896512X^6\notag\\
 &+71997696X^8+66311168X^{10}+20500480X^{12}+1720320X^{14},
 \label{eq:P28-example}\\
 \mathsf R_4(X)={}&245760X^6+2334720X^4+3955712X^2
                                    +\frac{37765120}{63}.
 \label{eq:R28-example}
\end{align}
These are $2^{14}W_4$ and $2^{14}R_4$, respectively. Then
\begin{equation}\label{eq:ordinary-28}
 \boxed{\displaystyle
 \sum_{n=0}^{\infty}
    \{\mathsf P_4(x_n)A_2(n)A_8(n)-\mathsf R_4(x_n)\}
       =-\frac{169984}{3}.}
\end{equation}
The associated exact coefficient matrix has determinant
$\det M_4=-83/483840$ and extraction weights $(1,21,35,7)$.
The numerical example is not used to infer the formula: every coefficient
above follows from the proved finite polynomials.

\paragraph{What this resolves and what remains open.}
The optimal consecutive-prefix conjecture is now proved for every
even pair with smaller order $2$, as well as the diagonal family
already settled in the source.  The determinant factor
$2-B_{q-2}$ explains why a uniform real-sign argument is misleading:
its sign can change as the even Bernoulli numbers grow.  The nonzero
factor is instead proved arithmetically.  The following sections extend the determinant method to $p=4,6$ and give
an exact smaller determinant for every even pair. Nonvanishing for all
unequal pairs $8\le p<q$ remains open.

% END sections/03_moment_rank.tex


% BEGIN sections/03b_general_hankel.tex
% Optional structural continuation; all formulas are exact.
% Reuses the Padé polynomials pi_n, constants c_n and h_n from moment_rank.tex.
\section{A fixed-size determinant for every pair of even orders}
\label{sec:general-even-reduction}

The preceding argument has a useful extension even where the final
nonvanishing problem remains unsolved.  For $p=2a$, $q=2b$, with
$1\leq a\leq b$, the optimal-prefix matrix is an explicitly scaled
rank-$a$ perturbation of a shifted Bernoulli Hankel matrix.  The Hankel
matrix is always nonsingular.  Moreover, the remaining determinant
has size $a$, independently of $q$.

For clarity, the general inherited moment notation is
\[
 D_{p,q}(s)=\sum_{n\ge0}A_p(n)A_q(n)(n+1/2)^{-s},
 \qquad \re s>3-p-q,
\]
continued meromorphically thereafter, and $D_r=D_{p,q}(-2r)$.
The same centered-parity argument proves regularity at these even
arguments. The forward and inverse Bernoulli transforms
\eqref{eq:forward-foundation}--\eqref{eq:inverse-foundation} remain valid,
now with $C_j$ the constant term at $N\to\infty$ of
\begin{align*}
 &N^{2j+1}A_p(N)A_q(N)
 +\sum_{k=1}^N k^{2j+1-q}A_p(k)\\
 &\hspace{15mm}
 +\sum_{k=1}^N k^{2j+1-p}A_q(k)
 +\sum_{k=1}^N k^{2j+1-p-q}.
\end{align*}
For $j\ge b$ the inherited formula is
$C_j=\sum_{k=0}^{b-1}a_{jk}Z_k$, with $Z_0=-1/2$ and
$Z_k=\zeta(2k)$ for $k\ge1$. Here the coefficients depend on
the indicated pair $(p,q)$. With $d=2j+1$ and
$e=d-p-q+2k+1$, they are zero for $e<0$, and otherwise
\begin{equation}\label{eq:general-source-entries}
 a_{jk}=B_e\left\{
 \frac{\binom{d-q+1}{p-2k}}{d-q+1}
 +\frac{\binom{d-p+1}{q-2k}}{d-p+1}\right\}.
\end{equation}
Each summand is omitted unless its lower binomial index lies between
$2$ and its upper index. These are the exact general formulas of
\cite[Section 7]{CanonicalFifth}; the finite summation argument of
Section~\ref{sec:tail-foundation} has the same form. In the formal
coordinate order $C_0,\ldots,C_{b-1},Z_0,\ldots,Z_{b-1}$, the
complete first-$q$ determinant is $\det M_{a,b}/(4b-1)!!$.
The new theorems below evaluate or reduce the remaining square minor.

Put
\[
 \mu_j=\begin{cases}B_{2j}/(2j)!,&j\geq0,\\0,&j<0,\end{cases}
 \qquad H_{a,b}=(\mu_{t+k-a+1})_{0\leq t,k<b}.
\]
The convention $\mu_j=0$ for $j<0$ is part of this definition; it is
not an analytic continuation of the formula for positive indices.

\begin{lemma}[Factorization of the general source matrix]
\label{lem:general-even-factor}
Let $M_{a,b}=(a_{b+t,k})_{0\leq t,k<b}$, with the source's convention
that a binomial summand is omitted when its lower index is less than
$2$.  Define
\begin{align*}
 R_{tt}&=(2b+2t-2a+1)!,& E_{kk}&=\frac1{(2b-2k)!},\\
 U_{tk}&=\mathbf1_{k<a}\,\mu_{t+k-a+1}
 \frac{(2t+1)!(2b-2k)!}
 {(2a-2k)!(2b+2t-2a+1)!}.
\end{align*}
Then
\[
 M_{a,b}=R(H_{a,b}+U)E,
 \qquad \operatorname{rank}U=a.
\]
\end{lemma}

\begin{proof}
In the source entry formula put $j=b+t$ and
$e=2(t+k-a+1)$.  When $e\geq0$, cancellation of the factorial in
$B_e$ gives
\[
 a_{b+t,k}=\mu_{t+k-a+1}\left\{
 \frac{(2b+2t-2a+1)!}{(2b-2k)!}
 +\mathbf1_{k<a}\frac{(2t+1)!}{(2a-2k)!}\right\}.
\]
When $e<0$, both sides vanish by the stated convention.  This proves
the factorization.  Only the first $a$ columns of $U$ are nonzero.
Their first $a$ rows form an antitriangular matrix with nonzero
antidiagonal, because $\mu_{t+k-a+1}=0$ for $t+k<a-1$ and
$\mu_0=1$.  Thus its rank is exactly $a$.
\end{proof}

\begin{theorem}[Nonsingularity and inertia of all the shifted matrices]
\label{thm:shifted-H-inertia}
The real symmetric matrix $H_{a,b}$ is nonsingular for every
$1\leq a\leq b$.  Its numbers of positive and negative eigenvalues
are, respectively,
\[
 \begin{cases}
 (b-h,h),&a=2h,\\
 (h+1,b-h-1),&a=2h+1.
 \end{cases}
\]
Consequently
\[
 \operatorname{sgn}\det H_{a,b}=
 \begin{cases}(-1)^h,&a=2h,\\
 (-1)^{b-h-1},&a=2h+1.
 \end{cases}
\]
\end{theorem}

\begin{proof}
Euler's even-zeta evaluation gives
\[
 \mu_r=(-1)^{r-1}\,2\sum_{m\geq1}x_m^r,
 \qquad r\geq1,\qquad x_m=(2\pi m)^{-2}.
\]
Conjugate $H_{a,b}$ by $\operatorname{diag}(1,-1,1,\ldots)$ and
call the result $G$.  For $t+k\geq a$ its entries are
$(-1)^a2\sum_{m\geq1}x_m^{t+k-a+1}$; on $t+k=a-1$ they are
$(-1)^{a-1}$, and below that antidiagonal they are zero.
Every moment matrix
$\bigl(2\sum_m x_m^{i+j+s}\bigr)_{0\leq i,j<d}$ with $s\geq1$
is positive definite: its quadratic form is
$2\sum_m x_m^s P(x_m)^2$, and a nonzero polynomial cannot vanish
at all the distinct positive points $x_m$.

First let $a=2h$.  Partition after the first $h$ indices:
\[
 G=\begin{pmatrix}0&B\\B^{\mathsf T}&P\end{pmatrix}.
\]
Here $P$ is positive definite.  The first $h$ columns of $B$ form
an antitriangular square matrix with antidiagonal $-1$, so $B$ has
row rank $h$.  The Schur complement
$-BP^{-1}B^{\mathsf T}$ is negative definite.  Congruence elimination
therefore gives $b-h$ positive and $h$ negative eigenvalues.

Now let $a=2h+1$ and partition after the first $h+1$ indices:
\[
 G=\begin{pmatrix}ee^{\mathsf T}&B\\
 B^{\mathsf T}&-P\end{pmatrix},
 \qquad e=(0,\ldots,0,1)^{\mathsf T}.
\]
The matrix $P$ is positive definite.  In the first $h$ rows of $B$,
its first $h$ columns form an antitriangular matrix with
antidiagonal $1$.  Hence
$ee^{\mathsf T}+BP^{-1}B^{\mathsf T}$ is positive definite:
if its quadratic form vanishes, first $e^{\mathsf T}v=0$ and then
the indicated minor forces the remaining $h$ components of $v$ to
vanish.  Congruence elimination gives $h+1$ positive and
$b-h-1$ negative eigenvalues.  The case $a=b=1$ is simply
$H_{1,1}=(1)$; empty blocks in the argument may otherwise be
omitted.  Diagonal sign conjugation preserves inertia.
\end{proof}

There is a more explicit reduction than a matrix determinant lemma
with an unevaluated $H_{a,b}^{-1}$.  Retain the monic polynomials
$\pi_n$ and positive numbers $h_n,c_n$ constructed in the preceding
Pad\'e calculation.  Write
\[
 \pi_n(x)=\sum_{t=0}^n p_{nt}x^t,
 \qquad
 p_{nt}=\frac{2(2n+2t+1)!(2n+1)!}
 {(4n+2)!(2n-2t+1)!(2t)!}.
\]
For $0\leq i,k<a$ set $n_i=b-a+i$, and define the $a$-square matrix
\begin{equation}\label{eq:general-small-A}
 A_{ik}=\sum_{t=0}^{n_i}p_{n_i t}\mu_{t+k-a+1}
 \left\{1+
 \frac{(2t+1)!(2b-2k)!}
 {(2a-2k)!(2b+2t-2a+1)!}\right\}.
\end{equation}
These are finite rational sums with explicitly specified coefficients.

\begin{theorem}[Reduction to a determinant of size $a$]
\label{thm:general-small-determinant}
For $1\leq a\leq b$,
\begin{equation}\label{eq:general-small-det}
 \det M_{a,b}=
 (-1)^{a(b-a)}
 \frac{\prod_{t=0}^{b-1}(2b+2t-2a+1)!}
      {\prod_{k=0}^{b-1}(2b-2k)!}
 \left(\prod_{n=0}^{b-a-1}h_n\right)\det A.
\end{equation}
All displayed scalar factors are nonzero.  Thus the remaining
optimal-prefix nonvanishing question is exactly $\det A\ne0$.
This reduction by itself does not assert that nonvanishing.
\end{theorem}

\begin{proof}
Let $\mathcal L_a(x^r)=\mu_{r-a+1}$ and recall
$\mathcal N(x^r)=\mu_{r+1}$.  Then
$\mathcal L_a(x^ag)=\mathcal N(g)$.  Replace row $n$ of
$H_{a,b}+U$ by the linear combination whose coefficients are those
of $\pi_n$, simultaneously for $0\leq n<b$.  This is a triangular
row transformation with diagonal $1$.

For a column $k\geq a$ the new entry is
\[
 \mathcal L_a(x^k\pi_n)=\mathcal N(x^{k-a}\pi_n).
\]
It vanishes when $k-a<n$, and equals $h_n$ when $k-a=n$.
The last $a$ rows consequently vanish in the last $b-a$ columns,
and their first $a$ columns are exactly $A$.  The first $b-a$ rows
in the last $b-a$ columns form an upper triangular matrix with
diagonal $h_0,\ldots,h_{b-a-1}$.  Moving the first $a$ columns past
the other $b-a$ columns gives the sign $(-1)^{a(b-a)}$.
Taking determinants and then restoring $R$ and $E$ proves the formula.
\end{proof}

For reference, the same reduction gives a fully explicit
small determinant for the unperturbed Hankel matrix.  Put
\[
 f_\ell(n)=\frac1{(2\ell+1)!}
 \prod_{r=0}^{2\ell-1}(2n-r)(2n+3+r),\qquad f_0(n)=1.
\]
Then
\begin{equation}\label{eq:general-H-small-det}
 \det H_{a,b}=(-1)^{a(b-a)}
 \left(\prod_{n=0}^{b-a-1}h_n\right)
 \left(\prod_{i=0}^{a-1}c_{n_i}\right)
 \det\bigl(f_{a-1-k}(n_i)\bigr)_{0\leq i,k<a}.
\end{equation}
Indeed,
\[
 \mathcal L_a(x^k\pi_n)
 =[z^{n+k-a+1}]F(z)q_n(z)
 =[z^{n+k-a+1}]r_n(z)
 =c_n f_{a-1-k}(n),
\]
where a negative coefficient index means zero.  The factorial
coefficient formula for $r_n$ proves the last equality, including
the zero when $n<a-1-k$.  In particular, the inertia theorem gives
the exact sign statement
\[
 (-1)^{a(a-1)/2}
 \det\bigl(f_{a-1-k}(b-a+i)\bigr)_{0\leq i,k<a}>0.
\]

\paragraph{A concrete route for the remaining even orders.}
For each fixed $a$, the entries in \eqref{eq:general-small-A} can be
reduced to finitely many nearby Bernoulli numbers by weighted
binomial Bernoulli identities.  More explicitly, if
$n=b-a+i$, $\ell=a-1-k$, and $N=2(n-\ell)+1\geq1$, the perturbation
part of $A_{ik}$ is
\[
 \frac{2(2n+1)!(2b-2k)!}
 {(4n+2)!(2a-2k)!N!}
 \sum_{r=0}^{n-\ell}\binom N{2r}B_{2r}
 (2r+2\ell+1)
 (2n+2\ell+2r+1)^{\underline{2i}}.
\]
The underline denotes a falling factorial.  The polynomial weight
has degree $2i+1\leq2a-1$.  Expressing it in the basis
$(N-2r)^{\underline d}$ and using
\[
 \sum_{j=0}^N\binom Nj B_j(N-j)^{\underline d}
   =N^{\underline d}B_{N-d}(1)
\]
therefore reduces it to Bernoulli numbers of orders at most $2a-1$
below $N$, with the omitted $j=1$ term restored explicitly.
When $n<\ell$ the perturbation entry is zero.  The small cases
$N=1,3,\ldots,2i+1$ require retaining $B_1(1)=1/2$ and $B_0=1$;
discarding them as generic odd-index zeros would be an error.
This is a finite algebraic program for each $p$, followed by a
genuine nonvanishing problem.  The sign theorem for $H_{a,b}$
does not determine the sign or nonvanishing of its perturbation.

% END sections/03b_general_hankel.tex


% BEGIN sections/03c_moment_rank_four.tex
% Follow-on research fragment. Requires the definitions and lemmas of
% moment_rank.tex, including pi_n, c_n, h_n, and Bernoulli 2-adic parity.

\section{The optimal prefix for every pair \texorpdfstring{$(4,q)$}{(4,q)}}
\label{rf:section}

The same orthogonal polynomials settle a second infinite unequal-order
family.  The perturbation now has rank two.  Its determinant involves
three consecutive even Bernoulli numbers, and their elementary
$2$-adic parity proves that this determinant never vanishes.

Let $p=4$, $q=2b$, and $b\geq3$.  In this section write
\[
 M_b^{[4]}=(a_{b+t,k})_{0\leq t,k<b}.
\]
For $b\geq4$, put
\[
 d_b=4b^2-6b-1,\qquad
 u=B_{2b-6},\quad v=B_{2b-4},\quad w=B_{2b-2},
\]
and define
\begin{align}
 \Xi_b={}&4d_b
 +(b-2)(2b-5)\{u+(2b-1)w(4-u)\}\notag\\
 &+(b-1)(2b-3)\{(2b-1)v^2-4(b+1)v\}.
 \label{rf:Xi}
\end{align}

\begin{theorem}[Exact determinant for the entire order-four family]
\label{rf:main}
For every integer $b\geq4$,
\begin{equation}\label{rf:det}
 \boxed{\displaystyle
 \det M_b^{[4]}=-\frac{\Xi_b}
      {6\,2^b(b-1)(2b-4)!}\ne0.}
\end{equation}
Moreover $4\Xi_b$ belongs to $\mathbb Z_{(2)}$ and is odd there;
equivalently, $v_2(\Xi_b)=-2$.  At the remaining unequal-order
case $b=3$,
\begin{equation}\label{rf:smallcase}
 M_3^{[4]}=
 \begin{pmatrix}
 0&3/4&1/4\\5/12&2/3&-1/12\\1&-3/8&1/12
 \end{pmatrix},\qquad
 \det M_3^{[4]}=-\frac{113}{384}.
\end{equation}
Thus the first $q$ centered moments have full formal rank $q$
for every even pair $(p,q)=(4,2b)$ with $q>4$.  Their complete
coefficient determinant is the determinant in
\eqref{rf:det} or \eqref{rf:smallcase}, divided by $(4b-1)!!$.
\end{theorem}

\subsection{A shifted Hankel matrix and its first two inverse rows}

Retain $\mu_n=B_{2n}/(2n)!$ for $n\geq0$, and put $\mu_{-1}=0$.
Define
\[
 H_b^{[-1]}=(\mu_{t+k-1})_{0\leq t,k<b}.
\]
The source formula gives the exact factorization
\begin{equation}\label{rf:factorization}
 M_b^{[4]}=\mathcal R_b(H_b^{[-1]}+U_b)\mathcal E_b,
\end{equation}
where
\[
 (\mathcal R_b)_{tt}=(2b+2t-3)!,\qquad
 (\mathcal E_b)_{kk}=\frac1{(2b-2k)!},
\]
and $U_b$ is zero outside its first two columns, with
\begin{equation}\label{rf:U}
 (U_b)_{tk}=\mu_{t+k-1}
  \frac{(2t+1)!(2b-2k)!}
       {(4-2k)!(2b+2t-3)!},\qquad k=0,1.
\end{equation}
This is precisely where the exclusion of a zero lower binomial
index matters: no third perturbation column is present.

Let $\mathcal L_{-1}(x^j)=\mu_{j-1}$, with
$\mathcal L_{-1}(1)=0$.  The shifted functional from the preceding
section satisfies
\[
 \mathcal L_{-1}(x^2g)=\mathcal N(g).
\]
Thus the explicit monic polynomials $\pi_n$ are orthogonal for the
functional on the right.  The Pad\'e coefficient identities already
proved give
\begin{align}
 \mathcal L_{-1}(x\pi_n)&=c_n,\notag\\
 \mathcal L_{-1}(\pi_n)&=c_n\mathfrak r_n,
 \qquad
 \mathfrak r_n=\frac{(2n-1)(2n)(2n+3)(2n+4)}6.
 \label{rf:two-moments}
\end{align}
The latter formula is valid also at $n=0$, where both sides vanish.
For $n\geq1$ it follows by taking the coefficient of $z^{n-1}$
in $F(z)q_n(z)$: that coefficient equals the coefficient of
$z^{n-1}$ in the explicit numerator $r_n(z)$.  The quotient of
this numerator coefficient and its leading coefficient $c_n$ is
exactly $\mathfrak r_n$.

Set
\[
 n=b-1,\quad m=b-2,\quad
 r_+=\mathfrak r_n,\quad r_-=\mathfrak r_m,
 \quad \delta_b=r_+-r_-
      =\frac{2(4b-3)d_b}{3}.
\]
Both $c_n$ and $c_m$ are nonzero, and $\delta_b>0$ for $b\geq2$.
Replace the last two monomial rows in $H_b^{[-1]}$ by the rows
for $\pi_m,\pi_n$.  This triangular row change has determinant
one, and these rows annihilate columns $2,\ldots,b-1$ by
$\mathcal N$-orthogonality.  Expanding by the first two columns
therefore gives
\begin{equation}\label{rf:Hdet}
 \det H_b^{[-1]}
 =-c_mc_n\delta_b\prod_{j=0}^{b-3}h_j\ne0.
\end{equation}
The two first rows of its inverse are consequently the coefficient
rows of the explicit polynomials
\begin{align}
 V_0(x)&=\frac1{\delta_b}
       \left(\frac{\pi_n(x)}{c_n}-\frac{\pi_m(x)}{c_m}\right),
       \label{rf:inverse0}\\
 V_1(x)&=\frac1{\delta_b}
       \left(-r_-\frac{\pi_n(x)}{c_n}
                       +r_+\frac{\pi_m(x)}{c_m}\right).
       \label{rf:inverse1}
\end{align}
Indeed, the two moments in \eqref{rf:two-moments} give the required
Kronecker values in columns zero and one, and orthogonality kills
all remaining columns.

Combining \eqref{rf:Hdet} with the diagonal factors gives
\begin{equation}\label{rf:base-scaled}
 \det\mathcal R_b\det\mathcal E_b\det H_b^{[-1]}
 =-\frac{d_b}{6\,2^b(b-1)(2b-4)!}\qquad(b\geq2).
\end{equation}
This factorial simplification can be checked inductively: after
division by $d_b$, the value at $b=2$ is $-1/24$, and the quotient
of successive values is $1/[4(b-1)(2b-5)]$.

\subsection{The four perturbation contractions}

For $j=n,m$ and $k=0,1$, put
\[
 s_{j,k}=\frac1{c_j}\sum_{t=0}^j[x^t]\pi_j(x)\,(U_b)_{tk},
 \qquad s_{+,k}=s_{n,k},\quad s_{-,k}=s_{m,k}.
\]
The explicit factorial coefficients of $\pi_j$ simplify these
contractions to the following finite Bernoulli sums:
\begin{align}
 s_{+,0}&=\frac{2b-2}{12}
 \sum_{\substack{0\leq k\leq2b-4\\k\ {\rm even}}}
 \binom{2b-3}{k}(k+3)(k+2b)(k+2b+1)B_k,\notag\\
 s_{+,1}&=\frac1{2b(2b-1)^2}
 \sum_{\substack{0\leq k\leq2b-2\\k\ {\rm even}}}
 \binom{2b-1}{k}(k+1)(k+2b-2)(k+2b-1)B_k,\notag\\
 s_{-,0}&=\frac{(2b)(2b-1)(2b-4)}{12}
 \sum_{\substack{0\leq k\leq2b-6\\k\ {\rm even}}}
 \binom{2b-5}{k}(k+3)B_k,\notag\\
 s_{-,1}&=\frac1{2b-3}
 \sum_{\substack{0\leq k\leq2b-4\\k\ {\rm even}}}
 \binom{2b-3}{k}(k+1)B_k.
 \label{rf:finite-sums}
\end{align}

For completeness, the polynomially weighted Bernoulli identities
needed to evaluate these sums are as follows.  If $N\geq3$ is odd,
\begin{equation}\label{rf:linear-sum}
 \sum_{k\ {\rm even}}\binom Nk(k+\alpha)B_k
 =N\left(\frac{\alpha+1}{2}-B_{N-1}\right).
\end{equation}
If $N\geq5$ is odd and $P$ is a monic cubic, then
\begin{align}
 \sum_{k\ {\rm even}}\binom NkP(k)B_k
 ={}&\frac N2P(1)
   +N\{P(N-1)-P(N)\}B_{N-1}\notag\\
 &-N(N-1)(N-2)B_{N-3}.
 \label{rf:cubic-sum}
\end{align}
Here and in \eqref{rf:linear-sum} the sum runs over $0\leq k\leq N$.
To prove the identities, first include the $k=1$ term, then expand
$P(N-y)$ in the falling factorials of $y=N-k$.  For $0\leq j\leq3$,
\[
 \sum_{k=0}^N(N-k)_{\underline j}\binom NkB_k
 =N_{\underline j}B_{N-j}(1).
\]
When $N\geq5$ is odd, only $j=1,3$ survive.  The falling-factorial
coefficient at $j=1$ is $P(N-1)-P(N)$, and the coefficient at
$j=3$ is $-1$.  Removing the $k=1$ term adds $NP(1)/2$.
The linear version is obtained in exactly the same way.  Thus these
are polynomial identities, not conjectural Bernoulli reductions.

Applying them to \eqref{rf:finite-sums}, for $b\geq4$ one obtains
\begin{align}
 s_{-,0}&=r_-(1-u/2),& s_{-,1}&=1-v,\notag\\
 s_{+,0}&=r_+
 -\frac{(2b-3)(2b-2)(2b-1)(4b-3)}3v\notag\\
 &\hspace{9mm}
 -\frac{(2b-3)(2b-2)(2b-4)(2b-5)}{12}u,\notag\\
 s_{+,1}&=1-\frac{2(4b-3)}b w
             -\frac{(b-1)(2b-3)}{b(2b-1)}v.
 \label{rf:four-contractions}
\end{align}
The separate treatment of $b=3$ avoids applying
\eqref{rf:linear-sum} at $N=1$, where that formula is not valid.

The rank-two determinant lemma and
\eqref{rf:inverse0}--\eqref{rf:inverse1} now give
\begin{equation}\label{rf:small-det}
 \frac{\det(H_b^{[-1]}+U_b)}{\det H_b^{[-1]}}
 =\frac{(r_++s_{+,0})(1+s_{-,1})
             -(r_-+s_{-,0})(1+s_{+,1})}{\delta_b}
 =\frac{\Xi_b}{d_b}.
\end{equation}
The last equality is direct polynomial simplification of the four
displayed contractions.  Together with \eqref{rf:base-scaled},
it proves the determinant formula \eqref{rf:det}.

\subsection{A uniform arithmetic proof of nonvanishing}

The Bernoulli parity argument in the order-two section proves that
\[
 \widehat u=2u,\qquad\widehat v=2v,\qquad\widehat w=2w
\]
are odd elements of $\mathbb Z_{(2)}$, since $b\geq4$ makes all
three Bernoulli indices positive and even.  Multiplication of
\eqref{rf:Xi} by four yields
\begin{align}
 4\Xi_b={}&16d_b
 +(b-2)(2b-5)
 \{2\widehat u+(2b-1)\widehat w(8-\widehat u)\}\notag\\
 &+(b-1)(2b-3)
 \{(2b-1)\widehat v^2-8(b+1)\widehat v\}.
 \label{rf:parity}
\end{align}
Modulo two, the first bracket contributes $b$ and the second
contributes $b-1$.  Hence
\[
 4\Xi_b\equiv b+(b-1)\equiv1\pmod2.
\]
Therefore $v_2(\Xi_b)=-2$ and $\Xi_b\ne0$ for every $b\geq4$.
This proves nonvanishing uniformly, without a real-sign assertion
or any assumption about Bernoulli numerators beyond the elementary
parity already established.

\paragraph{Updated scope of the initial-prefix problem.}
The optimal first-$q$ prefix is now proved for all even pairs with
$p=2$, all even pairs with $p=4$, and the diagonal $p=q$ inherited
from the source.  The next section settles $p=6$; the final remaining unequal-order
range in this article is $8\leq p<q$.  The rank-two calculation suggests investigating the
$2$-adic structure of the analogous higher-rank Bernoulli
contractions, but no general nonvanishing theorem for those
contractions is assumed here.

% Optional compact corollary of the explicit (4,6) source matrix.
\subsection{A convergent harmonic--polygamma identity at orders four and six}

The exceptional matrix \eqref{rf:smallcase} also gives a concrete
ordinary convergent identity.  Let $x_n=n+\tfrac12$ and
$A_s(n)=\zeta(s)-H_n^{(s)}$, and put
\[
 \begin{split}
 W_{4,6}(x)={}&67+15516x^2+324320x^4+1434496x^6\\
              &+1601280x^8+371712x^{10}.
 \end{split}
\]
Then
\begin{equation}\label{rf:ordinary-46}
 \boxed{\displaystyle
 \sum_{n=0}^{\infty}\left\{
 W_{4,6}(x_n)A_4(n)A_6(n)
 -\frac{123904x_n^2+316928}{5}\right\}=-21696.}
\end{equation}
The summand is $O(n^{-2})$, so the convergence is absolute.
One may equivalently write the tail product as
\[
 A_4(n)A_6(n)
 =\frac{\psi^{(3)}(n+1)\psi^{(5)}(n+1)}{720}.
\]

\begin{proof}
Direct multiplication of the exact matrix in
\eqref{rf:smallcase} gives
\[
 (-7,45,66)M_3^{[4]}=(339/4,0,0),
 \qquad -7C_3+45C_4+66C_5=-\frac{339}{8}.
\]
Set $S(y)=512(-7y^7+45y^9+66y^{11})$.  The inverse Bernoulli
transform identifies the polynomial of the resulting centered
moment combination as
\[
 S(x+\tfrac12)-S(x-\tfrac12)=W_{4,6}(x).
\]
The centered Hurwitz expansions give
\[
 A_4(n)A_6(n)=\frac1{15}x_n^{-8}
              -\frac7{60}x_n^{-10}+O(x_n^{-12}).
\]
Hence the polynomial in nonnegative powers of $x_n$ in its product
with $W_{4,6}$ is exactly
$(123904x_n^2+316928)/5$.  Subtracting it leaves an absolutely
summable remainder.  The removed polynomial has continued sum zero;
equivalently, its partial sum has zero constant term in the integer cutoff, by the
centered even-power argument used in
Corollary~\ref{cor:ordinary-two}.  The ordinary sum therefore equals
the finite-part moment combination, namely
$512(-339/8)=-21696$.
\end{proof}

% END sections/03c_moment_rank_four.tex


% BEGIN sections/03d_moment_rank_six.tex
% Exact p=6 continuation. Uses general_shifted_hankel.tex notation.
% Independent symbolic certificate: verify_moment_rank_six.py and .json.
\section{The optimal prefix for every pair \texorpdfstring{$(6,q)$}{(6,q)}}
\label{rs:section}

A congruence modulo four, rather than only the parity of doubled
Bernoulli numbers, settles the next unequal-order family.  All matrices
in this section are in the same source normalization as
Lemma~\ref{lem:general-even-factor}.

\begin{lemma}[A doubled-Bernoulli congruence]\label{rs:Bernoulli-mod-four}
For every integer $m\ge2$,
\[
 2B_{2m}\equiv1\pmod4\quad\hbox{in }\mathbb Z_{(2)}.
\]
\end{lemma}
\begin{proof}
The previously proved Bernoulli parity gives $v_2(B_{2j})=-1$
for $j\ge1$. Put $r=2m\ge4$. The exact Faulhaber identity is
\[
 \sum_{k=0}^3k^r
 =4B_r+\sum_{h=2}^{r+1}
          \binom r{h-1}\frac{4^h}{h}B_{r+1-h}.
\]
The $h=2$ term vanishes because $B_{r-1}=0$. Every $h\ge3$ term
is zero modulo eight: all Bernoulli numbers involved have valuation
at least $-1$, and
$2h-1-v_2(h)\ge3$. The left side is
$1+2^r+3^r\equiv2\pmod8$.
Thus $4B_r\equiv2\pmod8$, proving the assertion.
\end{proof}

\begin{theorem}[Order-six optimal prefix]\label{rs:main}
Let $a=3$ and $b\ge4$. Then
\[
 v_2\!\left(\frac{\det(H_{3,b}+U)}{\det H_{3,b}}\right)=-3.
\]
In particular $\det M_{3,b}\ne0$. Consequently the first $q$
centered moments have full formal rank $q$ for every pair
$(p,q)=(6,2b)$ with $q>6$.
\end{theorem}

Here are explicit finite identities certifying the uniform assertion.
They also specify the polynomial used in the proof without an
unrecorded fitted formula. For $0\le i,\ell\le2$, put
\[
 n=b-3+i,\qquad N=2(n-\ell)+1,\qquad R=2b-4+2\ell,
\]
and retain $f_\ell(n)$ from
\eqref{eq:general-H-small-det}. Define
\begin{align}
 F_{i\ell}&=f_\ell(n),\notag\\
 C_{i\ell}&=\frac{2R!}
 {(2n+1)(2n+2)(2\ell+2)!N!},\label{rs:C}\\
 W_{i\ell}(z)&=(z+2\ell+1)
                 (2n+2\ell+z+1)^{\underline{2i}},\notag\\
 \alpha_{i\ell h}&=\frac1{h!}\sum_{j=0}^h
             (-1)^{h-j}\binom hj W_{i\ell}(N-j).
 \label{rs:alpha}
\end{align}
All factorial ratios in $C_{i\ell}$ are interpreted as finite rational
products in $b$: $R-N=1+4\ell-2i$ is a fixed integer.
For independent variables $U_1,\ldots,U_5$, define the rational
$3\times3$ matrix
\begin{equation}\label{rs:matrix}
 \mathscr A_{i\ell}(b;U)
 =2F_{i\ell}+\frac{C_{i\ell}}2
 \sum_{\substack{1\le h\le2i+1\\h\ {\rm odd}}}
 \alpha_{i\ell h}N^{\underline h}
       U_{3-i+\ell+(h-1)/2}.
\end{equation}
Thus at most three terms occur in any entry. Put
\begin{align}
 D(b)&=16b^4-80b^3+80b^2+50b+9,\label{rs:D}\\
 Q(b)&=512b^6-6656b^5+27296b^4-25200b^3
              -91092b^2+217836b-123741.\label{rs:Q}
\end{align}
Finally define
\begin{equation}\label{rs:Pdefinition}
 \mathscr P(b;U)
 =48(2b-3)D(b)\,\frac{\det\mathscr A(b;U)}{\det F(b)}.
\end{equation}
The following are identities in rational functions of the indeterminates,
not assertions inferred from evaluations at integer values:
\begin{align}
 \det F(b)&=\frac{16}{45}(b-1)(2b-3)(4b-7)(4b-5)(4b-3)D(b),
 \label{rs:Fdet}\\
 \mathscr P(b;U)&\in\mathbb Z[b,U_1,\ldots,U_5],
 \label{rs:integrality}\\
 \mathscr P(b;1,1,1,1,1)&=-2(2b-5)Q(b).
 \label{rs:ones}
\end{align}
These finite algebraic identities are obtained by expanding the six
terms of each $3\times3$ determinant in
\eqref{rs:Pdefinition} and cancelling the displayed factors.
The accompanying certificate performs the expansion with an
indeterminate $b$, checks membership in the integer polynomial ring,
and verifies \eqref{rs:Fdet} and \eqref{rs:ones} exactly. It records
the full polynomial $\mathscr P$, which has 22 monomials in
$U_1,\ldots,U_5$ with integer-polynomial coefficients in $b$.
In particular \eqref{rs:matrix}--\eqref{rs:Pdefinition} are a finite
explicit definition from which every coefficient can be recovered.

\begin{proof}[Proof of Theorem~\ref{rs:main}]
Use the last three orthogonal-polynomial rows in
Theorem~\ref{thm:general-small-determinant}, reversing their first
three columns so that $\ell=0,1,2$. Divide row $i$ by $c_n$.
The unperturbed rows are $F_{i\ell}$.
For $b\ge6$, the perturbation contraction is
\[
 C_{i\ell}\sum_{\substack{0\le j\le N\\j\ {\rm even}}}
          \binom Nj W_{i\ell}(j)B_j.
\]
Expand $W_{i\ell}(N-y)$ in falling factorials; its coefficients
are \eqref{rs:alpha}. The weighted Bernoulli identity
\[
 \sum_{j=0}^N\binom Nj (N-j)^{\underline h}B_j
       =N^{\underline h}B_{N-h}(1)
\]
reduces the full sum. Here $N\ge2i+3$, so only odd $h$ occur,
and every surviving Bernoulli index is positive and even.
Removing the $j=1$ term contributes $NW_{i\ell}(1)/2$.
Direct factorial cancellation gives
\[
 C_{i\ell}\frac N2W_{i\ell}(1)=f_\ell(n).
\]
Thus the normalized sum of the unperturbed row and its perturbation
is precisely \eqref{rs:matrix}, at
$U_j=2B_{2b-2j}$. The row and column factors cancel in the quotient,
so \eqref{rs:Pdefinition} proves
\begin{equation}\label{rs:ratio}
 8\frac{\det(H_{3,b}+U)}{\det H_{3,b}}
 =\frac{\mathscr P(b;2B_{2b-2},2B_{2b-4},\ldots,2B_{2b-10})}
             {6(2b-3)D(b)}.
\end{equation}
For $b\ge7$, all five doubled Bernoulli numbers are $1$ modulo
four by Lemma~\ref{rs:Bernoulli-mod-four}.
The integer-coefficient assertion \eqref{rs:integrality} and
\eqref{rs:ones} therefore show that the numerator of
\eqref{rs:ratio} is $2$ modulo four: $2b-5$ is odd, and
$Q(b)\equiv1\pmod2$. Its valuation is exactly one.
The denominator also has valuation one, since
$D(b)\equiv1\pmod2$. Hence the quotient on the left of
\eqref{rs:ratio} is a $2$-adic unit, which gives the claimed
valuation $-3$.

The three remaining cases are exact source-matrix calculations:
\[
\begin{array}{c|c|c}
 b&\det(H_{3,b}+U)/\det H_{3,b}&\det M_{3,b}\\\hline
 4&7527/1240&2509/5120\\
 5&713039/70280&-713039/6451200\\
 6&2968099/1116360&2968099/2786918400.
\end{array}
\]
Each ratio has valuation $-3$ and each determinant is nonzero.
The general factorization and the previously proved nonsingularity
of $H_{3,b}$ now establish the full claim.
\end{proof}

\paragraph{A limit of the valuation pattern.}
The valuations $-1,-2,-3$ proved for $a=1,2,3$ (with $b\ge4$
in the $a=2$ case) do not justify a
uniform assertion $v_2(\det(H_{a,b}+U)/\det H_{a,b})=-a$.
Already at $a=4,b=9$, an exact calculation gives
\[
 \frac{\det(H_{4,9}+U)}{\det H_{4,9}}
 =-\frac{334092867653747}{85606130610},
 \qquad v_2=-1.
\]
This contradicts that proposed valuation pattern even in the range
$b\ge2a$. It does not refute the optimal-prefix conjecture; the
determinant in this example is nonzero. The unequal-order conjecture
remaining after the present results is $8\le p<q$.

% END sections/03d_moment_rank_six.tex


% BEGIN sections/04_polylog.tex
\section{A polylogarithmic boundary identity at every even order}
\label{sec:polylog-companion}

The new binomial annihilator has an exact polylogarithmic realization.
We use the one-argument, decreasing-index convention
\[
 \operatorname{Li}_{p,q}(z)=\sum_{n>m\ge1}\frac{z^n}{n^p m^q},
 \qquad |z|<1.
\]
Thus only the largest summation index carries the variable. This
convention is essential: a product of two ordinary polylogarithms
would give a different coloring.

\begin{proposition}[The two-tail generating function]
\label{prop:two-tail-polylog}
For integers $p,q\ge2$ and $|z|<1$, define
\[
 \mathcal F_{p,q}(z)=\sum_{n\ge0}
       \bigl(\zeta(p)-H_n^{(p)}\bigr)
       \bigl(\zeta(q)-H_n^{(q)}\bigr)z^n.
\]
Then
\begin{equation}\label{eq:two-tail-polylog}
 \mathcal F_{p,q}(z)=\frac{1}{1-z}\Bigl\{
 \zeta(p)\zeta(q)-\zeta(q)\operatorname{Li}_p(z)
 -\zeta(p)\operatorname{Li}_q(z)
 +\operatorname{Li}_{p,q}(z)+\operatorname{Li}_{q,p}(z)
 +\operatorname{Li}_{p+q}(z)\Bigr\}.
\end{equation}
\end{proposition}
\begin{proof}
Put $T_n=(\zeta(p)-H_n^{(p)})(\zeta(q)-H_n^{(q)})$.
The exact difference is
\[
 T_{n-1}-T_n=n^{-p}(\zeta(q)-H_n^{(q)})
            +n^{-q}(\zeta(p)-H_n^{(p)})+n^{-p-q}.
\]
Multiplication of the generating series by $1-z$ and substitution
of this difference give the result, using
\[
 \sum_{n\ge1}\frac{H_n^{(q)}}{n^p}z^n
   =\operatorname{Li}_{p,q}(z)+\operatorname{Li}_{p+q}(z).
\]
All series converge absolutely and locally uniformly on the open
unit disk, so the rearrangements are justified.
\end{proof}

The generating-function identity is a routine consequence of finite
harmonic summation. The additional content is the particular
all-order cancellation furnished by the determinant theorem.
Write
\[
 \mathcal E=z\frac{d}{dz}+\frac12.
\]
For the polynomial $W_b$ in \eqref{eq:W-polynomial-two}, let
$R_b$ be the nonnegative-power part of the centered asymptotic
Laurent series
\begin{equation}\label{eq:polylog-subtraction}
 W_b(X)\zeta(2,X+\tfrac12)\zeta(2b,X+\tfrac12).
\end{equation}
Only finitely many Bernoulli coefficients enter $R_b$; it is an
explicit even polynomial of degree $2b-2$.

\begin{theorem}[An all-order polylogarithmic boundary evaluation]
\label{thm:polylog-boundary}
For every integer $b\ge2$,
\begin{equation}\label{eq:polylog-boundary}
 \boxed{\displaystyle
 \lim_{z\uparrow1}\left\{
 W_b(\mathcal E)\mathcal F_{2,2b}(z)
       -R_b(\mathcal E)\frac{1}{1-z}\right\}
 =-\frac{2b-1}{4}\bigl(2-B_{2b-2}\bigr).}
\end{equation}
The expression in braces has an absolutely convergent Taylor series
on the closed unit disk, including its value at $z=1$.
\end{theorem}
\begin{proof}
On each compact subset of the open disk, polynomial differentiation
of the locally normally convergent generating series gives
\[
 W_b(\mathcal E)\mathcal F_{2,2b}(z)
       -R_b(\mathcal E)\frac{1}{1-z}
 =\sum_{n\ge0}\left\{
 W_b(n+\tfrac12)A_2(n)A_{2b}(n)-R_b(n+\tfrac12)
 \right\}z^n.
\]
The centered Hurwitz expansion contains only even inverse powers.
After the displayed subtraction its coefficients are $O(n^{-2})$.
The series therefore converges absolutely and uniformly for $|z|\le1$.
Evaluation at $z=1$ is exactly the ordinary convergent identity in
Corollary~\ref{cor:ordinary-two}.
\end{proof}

Equations~\eqref{eq:two-tail-polylog} and
\eqref{eq:polylog-boundary} give a finite prescription involving
classical polylogarithms, depth-two polylogarithms, and their ordinary
derivatives. The subtraction is determined before the limit is taken.
There is consequently no unspecified regularization constant and
no need to assign values to the divergent pieces separately.

% END sections/04_polylog.tex


% BEGIN sections/05_tornheim.tex
% Standalone section contributed to the 2026-10-11 research continuation.
% Required macros/environments: theorem, lemma, corollary, proposition, proof.
% Citation keys: TC:BD, TC:Romik, TC:Onodera2014, TC:Onodera2021.
\section{Third derivatives of Tornheim functions: exact cancellation}
\label{tc:sec}

Write
\[
 T(A,B,C)=\sum_{m,n\ge1}m^{-A}n^{-B}(m+n)^{-C},
 \qquad W_{a,b,c}(t)=T(at,bt,ct).
\]
The sum is initially absolutely convergent when
$\Re(A+C)>1$, $\Re(B+C)>1$, and $\Re(A+B+C)>2$.
Thereafter $T$ denotes its meromorphic continuation.  The normalization
contains no factor $2^t$ on the diagonal.  The origin is a crossing of
polar hyperplanes, so a derivative of $W_{a,b,c}$ is a derivative of
the specified one-variable restriction, not a partial derivative of
$T$ at the origin.

Borwein--Dilcher and Romik treat the diagonal and related first
derivatives; Onodera's work treats functional relations and first and
second derivative values in a broader Hurwitz setting
\cite{TC:BD,TC:Romik,TC:Onodera2014,TC:Onodera2021}.
The previous ProveIt continuation already supplies the entire
quadratic directional formula and a reduction of the cubic digamma
finite part to $W_{1,1,1}'''(0)$.  We do not repeat that reduction as a
new evaluation.  The result here describes every third directional
derivative with two specified constants, proves exactly when these
two coordinates cancel formally, and gives identities in ordinary
zeta derivatives alone.  This is a structural continuation of the
inspected material; no historical priority over all possible
specializations of earlier functional relations is asserted.

Put
\[
 p=\zeta'(0)\zeta''(0),\qquad q=\zeta'''(0),\qquad
 r=\zeta'''(-1),\qquad \Omega=W_{1,1,1}'''(0),
\]
and define
\begin{align}
 \Phi(a,b,c)&=c\{a^2+b^2-c(a+b)\},\label{tc:Phi}\\
 \mathscr Q(a,b,c)
 &=3(a+b)(ab+c^2)p
   -\left\{\frac{a^3+b^3}{2}+3c^2(a+b)+c^3\right\}q\notag\\
 &\quad+\left\{c^3-c\left(\frac{a^3}{b+c}
                              +\frac{b^3}{a+c}\right)\right\}r.
 \label{tc:Q}
\end{align}
All these functions are homogeneous of degree three on
$\mathscr D=\{(a,b,c):(a+c)(b+c)\ne0\}$.

\begin{theorem}[Complete third directional formula]
\label{tc:main}
There is a single direction-independent constant $\kappa$, defined
explicitly in \eqref{tc:kappa-def} below, such that every
$(a,b,c)\in\mathscr D$ has a holomorphic restriction at $t=0$ and
\begin{equation}
 \boxed{\displaystyle
 W_{a,b,c}'''(0)
   =abc\{\Omega-12p+8q\}
      +\Phi(a,b,c)\kappa+\mathscr Q(a,b,c).}
 \label{tc:third}
\end{equation}
Thus all third directional derivatives are in the span of two
specified constants over the displayed ordinary-zeta algebra.
The assertion is a spanning result.  It does not assert arithmetic
independence of the two constants or their irreducibility to a
smaller constant algebra.
\end{theorem}

\subsection{A holomorphic remainder with a fixed normalization}

The pole term needed for the theorem is
\begin{align}
 \mathscr P(A,B,C)
 ={}&-\frac{\Gamma(1-A)\zeta(B-1)}{\Gamma(C)(A+C)}
     -\frac{\Gamma(1-B)\zeta(A-1)}{\Gamma(C)(B+C)}.
 \label{tc:polar}
\end{align}

\begin{lemma}[Canonical local decomposition]
\label{tc:local}
There is a unique holomorphic germ $J(A,B,C)$, symmetric in $A,B$,
such that
\begin{equation}
 T(A,B,C)=\zeta(A)\zeta(B)+\mathscr P(A,B,C)+C J(A,B,C).
 \label{tc:decomposition}
\end{equation}
In particular
\begin{equation}
 \kappa=3J_{AA}(0,0,0)+\frac{\gamma^3}{12}
       +\frac{\gamma\zeta(2)}4+\frac{\zeta(3)}6
       -3\gamma\zeta''(-1)
 \label{tc:kappa-def}
\end{equation}
is a well-defined constant.
\end{lemma}

\begin{proof}
Use the classical Mellin representation
\[
 \Gamma(C)T(A,B,C)
   =\int_0^\infty x^{C-1}
       \operatorname{Li}_A(e^{-x})\operatorname{Li}_B(e^{-x})\,dx
\]
and Jonqui\`ere's locally normally convergent expansion
\[
 \operatorname{Li}_A(e^{-x})
   =\Gamma(1-A)x^{A-1}
     +\sum_{j\ge0}\frac{(-x)^j}{j!}\zeta(A-j),
 \qquad 0<x<2\pi,
\]
with $A$ in a sufficiently small disc about zero.  Subtract from the
product the six terms
\begin{align*}
 S(A,B;x)={}&\Gamma(1-A)\Gamma(1-B)x^{A+B-2}
 +\Gamma(1-A)\zeta(B)x^{A-1}
 +\Gamma(1-B)\zeta(A)x^{B-1}\\
 &-\Gamma(1-A)\zeta(B-1)x^A
 -\Gamma(1-B)\zeta(A-1)x^B+\zeta(A)\zeta(B).
\end{align*}
The remainder after multiplication by $x^{C-1}$ is normally
integrable on $(0,1)$ near the spectral origin; fixed parameter
derivatives add only powers of $\log x$.  The integral on
$(1,\infty)$ is normally convergent by exponential decay.  Integrate
the six subtracted terms explicitly.  The first three have
denominators nonzero at the origin; the last three give
$\zeta(A)\zeta(B)/C$ and the two polar terms in
\eqref{tc:polar}.  Thus for a holomorphic symmetric germ $H$,
\[
 T=\mathscr P+\frac{\zeta(A)\zeta(B)+C H(A,B,C)}{\Gamma(1+C)}.
\]
Since $\Gamma(1)=1$, the difference
$\zeta(A)\zeta(B)\{1/\Gamma(1+C)-1\}$ is divisible by $C$ as a
holomorphic germ.  Dividing that difference by $C$ and adding
$H/\Gamma(1+C)$ defines $J$.  Uniqueness follows first for $C\ne0$
and then by holomorphy.  The initially convergent Mellin calculation
and uniqueness of meromorphic continuation justify every equality.
\end{proof}

On a ray, the two pole terms become
\[
 -\frac{c}{a+c}\frac{\Gamma(1-at)}{\Gamma(1+ct)}\zeta(bt-1)
 -\frac{c}{b+c}\frac{\Gamma(1-bt)}{\Gamma(1+ct)}\zeta(at-1).
\]
These are holomorphic at zero for $(a,b,c)\in\mathscr D$,
including $c=0$, which proves the removability assertion in
Theorem~\ref{tc:main}.

\subsection{The two determined second derivatives of the remainder}

All derivatives of $J$ in this subsection are at $(0,0,0)$.
Let
\[
 U=J_{AA}=J_{BB},\quad V=J_{AB},\quad
 X=J_{AC}=J_{BC},\quad Y=J_{CC},\quad M=U+2X.
\]
Two elementary restrictions determine $M$ and $Y$ exactly.
Counting positive pairs with a prescribed sum gives
\[
 T(0,0,C)=\zeta(C-1)-\zeta(C).
\]
Consequently \eqref{tc:decomposition} yields
\[
 C J(0,0,C)=\zeta(C-1)-\zeta(C)
                   -\frac{1}{6\Gamma(1+C)}-\frac14.
\]
Taking the coefficient of $C^3$ gives
\begin{equation}
 Y=\frac{r-q}{3}
       -\frac{\gamma^3-3\gamma\zeta(2)+2\zeta(3)}{18}.
 \label{tc:Y}
\end{equation}

For the other restriction, rearrangement of an absolutely convergent
double sum, followed by continuation, proves the classical
quasi-shuffle identity
\begin{equation}
 T(A,0,C)+T(C,0,A)=\zeta(A)\zeta(C)-\zeta(A+C).
 \label{tc:stuffle}
\end{equation}
The polar terms in this sum have the holomorphic continuation
\begin{align}
 \mathscr P(A,0,C)+\mathscr P(C,0,A)
 ={}&\frac{\Gamma(1-A)\Gamma(1-C)}{12}
       \frac{\sin(\pi A)+\sin(\pi C)}{\pi(A+C)}\notag\\
 &-\frac{\zeta(A-1)}{\Gamma(1+C)}
  -\frac{\zeta(C-1)}{\Gamma(1+A)}.
 \label{tc:polar-symmetric}
\end{align}
The quotient in the first line is entire locally, since
$\sin(\pi A)+\sin(\pi C)
=2\sin(\pi(A+C)/2)\cos(\pi(A-C)/2)$.
The coefficient of $A^2C$ in $C J(A,0,C)+A J(C,0,A)$ is $M/2$.
The corresponding coefficient in the first line of
\eqref{tc:polar-symmetric} is
$\gamma(\gamma^2+\zeta(2))/24$.  Coefficient comparison in
\eqref{tc:stuffle} therefore gives
\begin{equation}
 M=p-q+\gamma\zeta''(-1)
      +(\gamma^2-\zeta(2))\zeta'(-1)
      -\frac{\gamma^3+\gamma\zeta(2)}{12}.
 \label{tc:M}
\end{equation}
This calculation uses the symmetric two-variable identity before
specialization; no ambiguous origin value of an individual term is
introduced.

\subsection{Proof of the complete cubic formula}

Let $B(a,b,c)$ be the third derivative at zero of
$\zeta(at)\zeta(bt)+\mathscr P(at,bt,ct)$.  Taylor expansion of
the holomorphic remainder gives exactly
\begin{align}
 W_{a,b,c}'''(0)
 &=B(a,b,c)+3c(a^2+b^2)U+6abcV
                  +6c^2(a+b)X+3c^3Y\notag\\
 &=B(a,b,c)+3\Phi(a,b,c)U+6abcV
                  +3c^2(a+b)M+3c^3Y.
 \label{tc:third-before}
\end{align}
For transparent coefficient checking put $z_j=\zeta^{(j)}(-1)$.
The ordinary Gamma expansion gives
\begin{align*}
 B(a,b,c)={}&-\frac{a^3+b^3}{2}q+3ab(a+b)p
       -c\left(\frac{b^3}{a+c}+\frac{a^3}{b+c}\right)r\\
 &-3c\gamma(a^2+b^2)z_2
       -3c\{\gamma^2(2ab+c(a+b))
                         +\zeta(2)(2ab-c(a+b))\}z_1\\
 &+\frac c{12}\Bigl[
       \gamma^3\{(a+c)^2+(b+c)^2\}
       +3\gamma\zeta(2)(a^2+b^2-2c^2)\\
 &\hspace{38mm}+2\zeta(3)(a^2+b^2-c(a+b)+2c^2)\Bigr].
\end{align*}
For example this follows directly by differentiating
\[
 \log\frac{\Gamma(1-at)}{\Gamma(1+ct)}
  =\gamma(a+c)t+\frac{\zeta(2)}2(a^2-c^2)t^2
                    +\frac{\zeta(3)}3(a^3+c^3)t^3+O(t^4).
\]
Substitute \eqref{tc:Y} and \eqref{tc:M} into
\eqref{tc:third-before}.  Collecting the coefficients of
$\Phi$ and $abc$ gives
\begin{equation}
 W_{a,b,c}'''(0)
 =\Phi(a,b,c)\kappa
   +abc\{6V-6(\gamma^2+\zeta(2))\zeta'(-1)\}
   +\mathscr Q(a,b,c).
 \label{tc:third-two}
\end{equation}
The diagonal has $\Phi(1,1,1)=0$ and
$\mathscr Q(1,1,1)=12p-8q$.  Hence its specialization identifies
the expression in braces as $\Omega-12p+8q$, proving
Theorem~\ref{tc:main}.

An optional independent spectral description of the second constant is
\begin{equation}
 \kappa=-\frac12W_{1,0,2}'''(0)
                +6p-\frac{41}{4}q+\frac72r.
 \label{tc:kappa-Euler}
\end{equation}
Here $T(t,0,2t)=Z(2t,t)$ in the strict Euler double-zeta convention.
Formula~\eqref{tc:kappa-Euler} is a bridge, not a claimed evaluation
of that derivative by ordinary zeta constants.

\subsection{All formal cancellations and an explicit exact identity}

\begin{corollary}[Finite combinations cancelling both cubic coordinates]
\label{tc:cancellation}
Let $\{(a_j,b_j,c_j)\}_{j=1}^N\subset\mathscr D$ and
$w_1,\ldots,w_N\in\mathbb C$.  If
\begin{equation}
 \sum_jw_j a_jb_jc_j=0,\qquad
 \sum_jw_j\Phi(a_j,b_j,c_j)=0,
 \label{tc:criterion}
\end{equation}
then
\begin{equation}
 \sum_jw_jW_{a_j,b_j,c_j}'''(0)
       =\sum_jw_j\mathscr Q(a_j,b_j,c_j).
 \label{tc:cancelled}
\end{equation}
Conditions~\eqref{tc:criterion} are necessary and sufficient for
formal cancellation of the two coordinates in
Theorem~\ref{tc:main}.  Necessity here concerns treating those
coordinates as independent symbols; it is not an arithmetic
independence theorem for their actual values.
\end{corollary}

The locus $\Phi=0$ supplies especially short two-term identities.
For $c\ne0$ it is $a^2+b^2=c(a+b)$.  Its asymmetric positive
integer point $(3,6,5)$ gives
\begin{equation}
 \boxed{\displaystyle
 W_{3,6,5}'''(0)-90W_{1,1,1}'''(0)
   =81p-\frac{403}{2}q-\frac{245}{11}r.}
 \label{tc:two-ray-example}
\end{equation}
The coefficient $90=3\cdot6\cdot5$ cancels the diagonal
coordinate, while $\Phi(3,6,5)=0$ cancels the other one.

For a fixed triple $(a,b,c)$, let
$\phi_c=\Phi(a,b,c)$, $\phi_a=\Phi(b,c,a)$,
$\phi_b=\Phi(c,a,b)$.  A convenient family of weights for the three
cyclic rays is
\[
 (w_c,w_a,w_b)=(\phi_a-\phi_b,\phi_b-\phi_c,\phi_c-\phi_a).
\]
Both conditions in \eqref{tc:criterion} hold identically.
At $(a,b,c)=(1,2,3)$ these weights, divided by four, are $(1,4,-5)$.
Thus we obtain the concrete ordinary-zeta identity
\begin{equation}
 \boxed{\begin{aligned}
 &W_{1,2,3}'''(0)+4W_{2,3,1}'''(0)-5W_{3,1,2}'''(0)\\
 &\qquad=99\zeta'(0)\zeta''(0)
       +\frac{207}{2}\zeta'''(0)
       +\frac{162}{5}\zeta'''(-1).
 \end{aligned}}
 \label{tc:example}
\end{equation}
Every ray in this example has positive coordinates, so no zero-order
factor or restriction to a singular hyperplane is involved.
Equation~\eqref{tc:example} proves a new explicit cancellation relative
to the inspected ProveIt material; it does not evaluate $\Omega$.

\subsection{The exact ambiguity of these local identities at every order}

The appearance of two coordinates at third order has an algebraic
explanation that extends to all orders.  This is a statement about
what the preceding local restrictions determine, not about how many
arithmetically independent constants the actual Tornheim function has.

\begin{theorem}[Holomorphic ambiguity kernel]
\label{tc:kernel}
Suppose two meromorphic germs $T_1,T_2$ have the same polar part
\eqref{tc:polar}, in the precise sense that $T_i-\mathscr P$ is
holomorphic for $i=1,2$. Suppose both germs satisfy symmetry in their first two variables, have
the restriction $T(A,B,0)=\zeta(A)\zeta(B)$ away from the polar set,
and satisfy \eqref{tc:stuffle}.  Their difference has a unique form
\begin{align}
 \Delta T(A,B,C)={}&ABC\,S(A,B,C)\notag\\
 &+AC(A-C)R(A,C)+BC(B-C)R(B,C),
 \label{tc:kernel-form}
\end{align}
where $S$ is a holomorphic germ symmetric in $A,B$ and $R$ is a
holomorphic germ symmetric in its two variables.  Conversely every
expression in \eqref{tc:kernel-form} preserves these local identities.

For homogeneous degree $d\ge3$, the formal ambiguity space has
dimension
\begin{equation}
 \left\lfloor\frac{d^2-1}{4}\right\rfloor.
 \label{tc:kernel-dimension}
\end{equation}
It vanishes in degrees zero, one and two.  Fixing the entire diagonal
germ imposes $S(t,t,t)=0$, and reduces the degree-$d$ dimension by one.
\end{theorem}

\begin{proof}
The common polar term cancels.  The common restriction at $C=0$
then implies $\Delta T=C K(A,B,C)$ for a holomorphic germ $K$
symmetric in $A,B$.  Put $u(A,C)=K(A,0,C)$.  The difference of
\eqref{tc:stuffle} says
\[
 C u(A,C)+A u(C,A)=0.
\]
At $A=0$ this gives $u(0,C)=0$, so $u(A,C)=A v(A,C)$.
After division by $AC$ the identity becomes
$v(A,C)=-v(C,A)$.  An antisymmetric holomorphic germ vanishes on
$A=C$ and is uniquely divisible by $A-C$; its quotient is symmetric.
Thus $u(A,C)=A(A-C)R(A,C)$ for a unique symmetric germ $R$.
The remainder
\[
 K(A,B,C)-u(A,C)-u(B,C)
\]
vanishes at both $A=0$ and $B=0$.  It is therefore uniquely divisible
by $AB$, with a holomorphic quotient symmetric in $A,B$.  This
quotient is $S$ and proves \eqref{tc:kernel-form}.  Direct substitution
proves the converse and also uniqueness.

In homogeneous degree $d$, $S$ and $R$ have degree $d-3$.
Writing $m=d-3$, a symmetric polynomial in two selected variables
and a third unrestricted variable has
$\lfloor(m+2)^2/4\rfloor$ coefficients; a symmetric polynomial in
two variables of total degree $m$ has $\lfloor m/2\rfloor+1$
coefficients.  Their sum is $\lfloor(d^2-1)/4\rfloor$.
On the diagonal the $R$ terms vanish and
$\Delta T(t,t,t)=t^3S(t,t,t)$.  At each homogeneous degree evaluation
of $S$ at $(1,1,1)$ is a nonzero linear functional, so it removes
exactly one dimension.
\end{proof}

At $d=3$ the two free polynomials in \eqref{tc:kernel-form} are
constants; the two direction functions are exactly $abc$ and
$\Phi(a,b,c)$.  This proves that the cancellation criterion is the
complete one obtainable from these local identities alone.  At
degrees $4,5,6$ the corresponding formal dimensions are $3,6,8$.
Additional functional relations may reduce the arithmetic data
further; this theorem does not prejudge that question.

\paragraph{Verification and next questions.}
The accompanying exact symbolic calculation checks the coefficient
collection, the three-ray identity, and the formal kernel dimensions
independently by linear algebra.  A separate numerical script computes
the ray restrictions from the convergent Jonqui\`ere--Mellin formula
and an incomplete-Gamma tail.  Its floating-point residuals are
diagnostics and are not used in the proofs.  At 50 working decimal
digits, the independent comparison for \eqref{tc:example} gives
an absolute residual below $6.3\times10^{-33}$; both sides are
approximately $-451.3535195684862560644674145675405$.
Concrete next questions
are to identify additional functional relations constraining $R$,
to compute the full fourth-order formula in a similarly normalized
basis, and to determine whether the two actual cubic constants admit
arithmetic relations beyond the formal cancellation proved here.

% END sections/05_tornheim.tex


% BEGIN sections/06_gamma_foundation.tex
\section{Height-one Stieltjes coefficients and the inherited Gamma bridge}
\label{sec:gamma-foundation}

The sixth and seventh reductions use a height-one Laurent coefficient.
We record its analytic definition and prove the inherited integral bridge
in the exact convention needed here. These coefficient formulas are
established in the incoming fifth-moment report and are related to the
height-one series of Young \cite{CanonicalFifth,Young2024}; their inclusion
does not constitute a new evaluation of generalized Stieltjes constants.

For depth $d\ge1$, let
\[
 Z_d(s)=\zeta(s,\underbrace{1,\ldots,1}_{d-1})
       =\sum_{n_1>\cdots>n_d\ge1}
          \frac{1}{n_1^s n_2\cdots n_d},\qquad \re s>1.
\]
Define $c_j$ and $\lambda_{d,j}$ by
\begin{align}
 \frac1{\Gamma(1+u)}&=\sum_{j\ge0}c_ju^j,\label{eq:gamma-c-foundation}\\
 Z_d(1+u)&=\sum_{j=0}^d c_ju^{j-d}
                    +\sum_{j\ge1}\lambda_{d,j}u^j,
 \qquad \lambda_d:=\lambda_{d,1}.
 \label{eq:height-one-foundation}
\end{align}
The second expansion will follow from the next proposition. In particular,
$\lambda_{1,j}=(-1)^j\gamma_j/j!$ and $\lambda_1=-\gamma_1$.
The multiple coefficients $\lambda_{d,j}$ have a specified depth; they
should not be identified with ordinary Hurwitz coefficients $\gamma_j(a)$
without a further identity.

For $v>0$, set
\[
 B(v)=-\log(1-e^{-v}),\qquad
 V_d(v)=B(v)^d-\boldsymbol1_{(0,1)}(v)(-\log v)^d,
\]
and put
\[
 I_{d,r}=\int_0^\infty\frac{(\log v)^rV_d(v)}v\,dv
 \quad(r\ge0).
\]
At zero, $V_d(v)=O(v(1+|\log v|^{d-1}))$; at infinity it
decays exponentially. All these integrals converge absolutely.

\begin{proposition}[Fixed Mellin normalization]
\label{prop:height-one-foundation}
For $d\ge1$, as a meromorphic identity in $\re u>-1$,
\begin{equation}\label{eq:height-one-Mellin-foundation}
 \Gamma(1+u)Z_d(1+u)
 =u^{-d}+\frac{u}{d!}
          \int_0^\infty v^{u-1}V_d(v)\,dv.
\end{equation}
Consequently the coefficients in \eqref{eq:height-one-foundation} satisfy
\begin{equation}\label{eq:height-one-allcoeff-foundation}
 \lambda_{d,j}=c_{d+j}
       +\frac1{d!}\sum_{r=0}^{j-1}
                     \frac{c_{j-1-r}}{r!}I_{d,r},\qquad j\ge1.
\end{equation}
If $\Lambda_d(u)=\Gamma(1+u)Z_d(1+u)-u^{-d}$, then
$\Lambda_d^{(j)}(0)=jI_{d,j-1}/d!$ for $j\ge1$.
\end{proposition}

\begin{proof}
The strict-index generating series with $d-1$ unit indices is
$(-\log(1-y))^{d-1}/(d-1)!$. Termwise Mellin integration for
$\re u>0$ therefore gives
\[
 \Gamma(1+u)Z_d(1+u)
  =\frac1{(d-1)!}\int_0^\infty
            \frac{v^u B(v)^{d-1}}{e^v-1}\,dv
  =\frac{u}{d!}\int_0^\infty v^{u-1}B(v)^d\,dv.
\]
For the second equality integrate by parts, using
$B'(v)=-(e^v-1)^{-1}$; both boundary terms vanish in this half-plane.
Subtract the exact integral
$\int_0^1v^{u-1}(-\log v)^d\,dv=d!/u^{d+1}$.
The estimates above give holomorphy of the remaining integral for
$\re u>-1$, locally uniformly with every spectral derivative.
Expansion at zero and multiplication by \eqref{eq:gamma-c-foundation}
prove all the coefficient identities.
\end{proof}

This calculation makes the Stieltjes content of the later formulas
explicit: the spectral derivative data are ordinary convergent logarithmic
Mellin integrals, with the subtraction at $v=1$ fixed in advance.

Define the Gamma family
\begin{equation}\label{eq:Cgamma-foundation}
 G(x,t)=\frac{\Gamma(1-t)\Gamma(x)}{\Gamma(x-t)}
       =\sum_{m\ge0}C_m(x)\frac{t^m}{m!}.
\end{equation}
At positive integer $x=n$, it is the familiar complete harmonic-star
generator $\prod_{k=1}^{n-1}(1-t/k)^{-1}$; this classical connection
is also discussed by Sesma \cite{Sesma}. At $x=0$, the residue of
$G(x,t)$ is $-t$. Thus every $C_m$ with $m\ge2$ has a removable
singularity there. Expansion of the constant term gives
$C_m(0)=-m!\zeta(m)$, while $G(1,t)=1$ gives $C_m(1)=0$.

\begin{proposition}[Inherited ordinary Gamma integral]
\label{prop:Cgamma-integral-foundation}
For every $m\ge2$,
\begin{equation}\label{eq:Cgamma-integral-foundation}
 \int_0^1C_m(x)\,dx=m!(c_m-\lambda_{m-1}).
\end{equation}
The integral is ordinary and the complete integrand is continuous on
$[0,1]$.
\end{proposition}

\begin{proof}
Use $F(x,t)=G(x,-t)$ and write $h_1(x)=\psi(x)+\gamma$. The function
\[
 \Psi(t)=\int_0^1\{F(x,t)-1+t h_1(x)\}\,dx
\]
is holomorphic near zero. Indeed the residues $t/x$ and $-t/x$
cancel, leaving a jointly regular expression for $x\in[0,1]$ and
$t$ in a sufficiently small closed disc. Hence
\[
 \Psi(t)=\sum_{m\ge2}\frac{(-1)^m t^m}{m!}
                                  \int_0^1C_m(x)\,dx.
\]
For $\re t>0$, Euler's beta integral for $F$ and the ordinary
digamma integral, combined before integration, yield
\[
 \Psi(t)=t\int_0^\infty
   \left\{\frac{(1-e^{-v})^t-1}{v}
                    +\frac{e^{-v}}{1-e^{-v}}\right\}\,dv-1.
\]
Fubini is justified for the combined integrand: at zero it is
$O(v^{\re t-1}+1)$ after cancellation, and at infinity integration
in $x$ leaves an exponentially decaying numerator divided by $v$.
To put this in the fixed subtraction convention, integrate from a
lower cutoff $\delta$, use
$\int_\delta^\infty(e^v-1)^{-1}\,dv=-\log(1-e^{-\delta})$,
and compare with $\int_\delta^1(v^t-1)\,dv/v$. Taking
$\delta\downarrow0$ gives
\begin{equation}\label{eq:gamma-Psi-foundation}
 \Psi(t)=t\int_0^\infty
 \frac{(1-e^{-v})^t-1
       -\boldsymbol1_{(0,1)}(v)(v^t-1)}v\,dv.
\end{equation}
Near zero, $(1-e^{-v})^t-v^t=O(v^{\re t+1})$ uniformly on compact
parameter sets; at infinity the numerator decays exponentially.
Thus this equality continues holomorphically to a neighborhood of zero,
and coefficients may be extracted under the integral. Its degree-$m$
coefficient is $(-1)^{m-1}I_{m-1,0}/(m-1)!$.
Comparing the two expressions for $\Psi$ gives
$\int_0^1C_m=-mI_{m-1,0}$. Proposition~\ref{prop:height-one-foundation}
then proves \eqref{eq:Cgamma-integral-foundation}.
\end{proof}

Individual terms in a differential expansion of $C_m$ can have divergent
improper integrals even though $C_m$ itself is regular. The finite-part
normal form in the next section is what permits those terms to be
separated with their exact endpoint contributions retained.

% END sections/06_gamma_foundation.tex


% BEGIN sections/06_differential.tex
% Research component for root integration. All labels have the dq prefix.
% Bibliography required: OlverShakiban1978; incoming canonical fifth report.
% This is a quotient of VECTOR SPACES, not a quotient algebra by an ideal.

\section{Canonical integration by parts for polygamma polynomials}
\label{dq:section}

The incoming canonical fifth-moment report asks for an all-degree
classification of differential monomials after integration by parts,
with their finite endpoint constants retained. The next theorem answers
the formal integration-by-parts part of that question. It also gives an
exact reduction algorithm and a smaller parity sector for the inverse-Gamma
polynomials. The vector-space dimensions below do not count arithmetically
independent constants, or independent images under any particular spectral
map. Additional special-function identities can reduce the resulting
integrals further; the sixth-degree calculation below does exactly this.

Differential-polynomial methods and the exactness of the Euler-operator
complex are established tools; see Olver--Shakiban
\cite{OlverShakiban1978}. The elementary construction below is included
in full to fix the particular grading, basis, and endpoint convention
needed for this application.

\subsection{The formal vector-space quotient and a triangular basis}

Let
\[
 \mathscr R=\mathbb Q[X_1,X_2,\ldots],\qquad
 D=\sum_{j\ge1}X_{j+1}\frac{\partial}{\partial X_j}.
\]
Each polynomial involves finitely many variables, so this derivation
is well defined. Give $X_j$ weight $j$ and factor degree one. Write
$\mathscr R_{N,\ell}$ for the span of monomials of weight $N$ and
factor degree $\ell$, and $\mathscr R_N=\bigoplus_\ell\mathscr R_{N,\ell}$.
The derivation raises weight by one and preserves factor degree.
The substitution $X_j=\psi^{(j-1)}(x)$ intertwines $D$ with $d/dx$.

For a partition $\lambda=(\lambda_1\ge\cdots\ge\lambda_\ell>0)$,
put $X_\lambda=\prod_{i=1}^\ell X_{\lambda_i}$.
Let $\mathscr B_N$ consist of the partitions of $N$ whose largest
part occurs at least twice. At weights zero and one, use separately
the representatives $1$ and $X_1$.

\begin{theorem}[Canonical differential-polynomial normal form]
\label{dq:normalform}
For $N\ge2$, the classes
\[
 \{X_\lambda:\lambda\in\mathscr B_N\}
\]
are a basis of the vector space $\mathscr R_N/D\mathscr R_{N-1}$.
Every $F\in\mathscr R_N$ has a unique decomposition
\begin{equation}\label{dq:decomposition}
 F=\mathcal N(F)+D\mathcal T(F),\qquad
 \mathcal N(F)\in\operatorname{span}_{\mathbb Q}
                 \{X_\lambda:\lambda\in\mathscr B_N\},
 \quad\mathcal T(F)\in\mathscr R_{N-1}.
\end{equation}
If $p_\ell(N)$ is the number of partitions of $N$ into exactly
$\ell$ parts, then
\begin{equation}\label{dq:dimensions}
 \dim \frac{\mathscr R_{N,\ell}}{D\mathscr R_{N-1,\ell}}
   =p_\ell(N)-p_\ell(N-1),\qquad
 \dim \frac{\mathscr R_N}{D\mathscr R_{N-1}}
   =p(N)-p(N-1).
\end{equation}
The decomposition is produced by a terminating rational algorithm.
Every retained monomial uses derivatives of order at most
$\lfloor N/2\rfloor-1$ after substitution into $\psi$.
\end{theorem}

\begin{proof}
Order partitions lexicographically by their decreasing lists of parts.
Let a monomial have a unique largest part $M\ge2$. Decrease this
part to $M-1$, and call the resulting partition $\mu$. If $b$ is
the multiplicity of $M-1$ in $\mu$, the product rule gives
\begin{equation}\label{dq:pivot}
 D X_\mu=bX_\lambda+R_\lambda.
\end{equation}
Every monomial in $R_\lambda$ has largest part at most $M-1$:
it arises by differentiating one of the smaller factors of $X_\mu$.
Thus every such partition is lexicographically smaller than $\lambda$.
Replace $X_\lambda$ by $b^{-1}DX_\mu-b^{-1}R_\lambda$.
Applying this replacement to the largest remaining unique-largest
partition strictly decreases a finite ordered set of partitions of $N$,
so it terminates in the asserted remainder. The replacement preserves
factor degree.

For independence, apply $D$ to the partition monomials of weight $N-1$.
The leading partition is obtained by increasing one largest part by one;
its coefficient is the multiplicity of that largest part and is nonzero.
This is a bijection between partitions of $N-1$ and partitions of $N$
with a unique largest part. It also preserves the number of parts.
Distinct derivative columns therefore have distinct leading partitions.
No nonzero linear combination of those columns can be supported entirely
on repeated-largest partitions. This proves both independence and
injectivity of $D$ on $\mathscr R_{N-1}$, and hence uniqueness of
the primitive in \eqref{dq:decomposition}. Counting the columns and
rows gives \eqref{dq:dimensions}. Finally $2\lambda_1\le N$ for
each retained partition, giving the derivative-order bound.
\end{proof}

The quotient is a vector-space quotient: $D\mathscr R$ is not in
general an ideal, so multiplication of classes is not being asserted.
The ordinary generating series, including weights zero and one, is
\begin{equation}\label{dq:Hilbert}
 \sum_{N\ge0}\dim(\mathscr R_N/D\mathscr R_{N-1})q^N
 =q+\prod_{j\ge2}\frac1{1-q^j}.
\end{equation}
The convention is $D\mathscr R_{-1}=0$. The extra $q$ records
$X_1$, since differentiation kills the constant polynomial.
With factor degree recorded by $z$, the corresponding series is
\[
 (1-q)\prod_{j\ge1}(1-zq^j)^{-1}+q.
\]

For example, the respective bases at weights four, five, and six are
\[
 \begin{array}{c|l}
 4& X_1^4,\ X_2^2\\
 5& X_1^5,\ X_1X_2^2\\
 6& X_1^6,\ X_1^2X_2^2,\ X_2^3,\ X_3^2.
 \end{array}
\]
This explains formally why the fifth digamma calculation introduces
$\psi(\psi')^2$ after the scalar power. It does not assert that its
integral is independent of the spectral constants already present.

\subsection{The exact endpoint functional}

For a meromorphic function with a finite principal part at zero, use
the original-coordinate finite part
\[
 \operatorname{FP}\int_0^1f(x)\,dx
 =\operatorname{CT}_{\varepsilon\downarrow0}\int_\varepsilon^1 f(x)\,dx.
\]
Here the constant term is taken after all negative powers of
$\varepsilon$ and all nonconstant powers of $\log\varepsilon$ are removed.
For every differential polynomial $T$,
\begin{equation}\label{dq:endpointfunctional}
 \operatorname{FP}\int_0^1 DT(\psi,\psi',\ldots)\,dx
 =\mathcal E(T):=T(\psi(1),\psi'(1),\ldots)
          -\operatorname{CT}_{x=0}T(\psi(x),\psi'(x),\ldots).
\end{equation}
This follows by the fundamental theorem of calculus before the constant
term is taken. In particular a total derivative cannot simply be
discarded in the finite-part problem.

The functional is an explicitly finite polynomial calculation. At
weight $w$, use only
\begin{equation}\label{dq:endpointtruncation}
 U_w(x)=-x^{-1}-\gamma+
       \sum_{k=1}^{w-1}(-1)^{k+1}\zeta(k+1)x^k,
\end{equation}
and substitute $X_j=U_w^{(j-1)}(x)$ to obtain the constant term.
Terms omitted from $\psi$ start at $x^w$; in a weight-$w$ monomial,
any occurrence of an omitted term has positive total exponent even
when all other factors contribute their leading poles. Thus it cannot
affect the constant coefficient. The endpoint at one uses
\[
 X_1=-\gamma,\qquad X_j=(-1)^j(j-1)!\zeta(j)\quad(j\ge2).
\]
Consequently $\mathcal E(T)\in\mathbb Q[\gamma,\zeta(2),\ldots,\zeta(w)]$
for every $T\in\mathscr R_w$.

For a formula with entirely finite summation ranges, define
$\beta_1=-\gamma$ and $\beta_j=(-1)^j(j-1)!\zeta(j)$ for $j\ge2$.
Set
\[
 \alpha_{j,-j}=(-1)^j(j-1)!,\quad \alpha_{1,0}=-\gamma,
 \quad
 \alpha_{j,k}=(-1)^{j+k}\frac{(j+k-1)!}{k!}\zeta(j+k)
 \quad(k\ge0,\ j+k\ge2).
\]
All other negative-index coefficients are zero. For a partition
$\lambda$ of $w$,
\begin{equation}\label{dq:finiteendpoint}
 \mathcal E(X_\lambda)=\prod_i\beta_{\lambda_i}
 -\sum_{\substack{
 k_i\in\{-\lambda_i,0,1,\ldots,w-\lambda_i\}\\
 k_1+\cdots+k_\ell=0}}
       \prod_i\alpha_{\lambda_i,k_i}.
\end{equation}
This supplies every endpoint constant at arbitrary degree without
leaving an unspecified finite-part operation.

The normal-form theorem therefore becomes the exact analytic identity
\begin{equation}\label{dq:integratednormalform}
 \operatorname{FP}\int_0^1F(\psi,\psi',\ldots)\,dx
 =\sum_{\lambda\in\mathscr B_N}
   [X_\lambda]\mathcal N(F)\,
   \operatorname{FP}\int_0^1\prod_i\psi^{(\lambda_i-1)}(x)\,dx
   +\mathcal E(\mathcal T(F)).
\end{equation}
Its proof uses finite algebra and \eqref{dq:endpointfunctional}; it
does not interchange divergent integrals of individual terms.

\subsection{Inverse-Gamma polynomials and the parity reduction}

Define formal differential polynomials by
\begin{equation}\label{dq:Arecurrence}
 A_0=1,\qquad A_{m+1}=X_1A_m-DA_m.
\end{equation}
On substituting $X_j=\psi^{(j-1)}$, these become
\[
 A_m(x)=(-1)^m\Gamma(x)\frac{d^m}{dx^m}\frac1{\Gamma(x)},
 \qquad
 \frac{\Gamma(x)}{\Gamma(x-t)}=
                 \sum_{m\ge0}A_m(x)\frac{t^m}{m!}.
\]
Write $N_m=\mathcal N(A_m)$ and $T_m=\mathcal T(A_m)$, with
$N_0=1,T_0=0,N_1=X_1,T_1=0$.

\begin{theorem}[Factor parity and an explicit all-degree recurrence]
\label{dq:parity}
Every monomial of $N_m$ has factor degree congruent to $m$ modulo two.
Thus its number of possible canonical coordinates is at most
\begin{equation}\label{dq:paritycount}
 r_m=\sum_{\ell\equiv m\ (2)}
             \bigl(p_\ell(m)-p_\ell(m-1)\bigr)\qquad(m\ge2).
\end{equation}
For $m\ge1$, the complete pair $(N_{m+1},T_{m+1})$ is computed
from $(N_m,T_m)$ by setting
\[
 (R,S)=(\mathcal N,\mathcal T)(X_1N_m-X_2T_m)
\]
and then
\begin{equation}\label{dq:pairrecurrence}
 N_{m+1}=R,\qquad
 T_{m+1}=S+X_1T_m-N_m-DT_m.
\end{equation}
All coefficient operations are finite and rational.
\end{theorem}

\begin{proof}
Here is a polynomial form of the adjoint argument, retaining its
boundary explicitly. Introduce a function $F$ with $F'=X_1$, and
write $p=e^F,q=e^{-F}$. The products below are differential
polynomials because the exponential factors cancel. The Lagrange
concomitant is
\[
 L_m=\sum_{j=0}^{m-1}(-1)^j p^{(j)}q^{(m-1-j)}.
\]
Termwise differentiation telescopes to
\[
 DL_m=pq^{(m)}+(-1)^{m-1}p^{(m)}q
       =(-1)^m A_m(X)-A_m(-X).
\]
Equivalently, put $B_j(X)=(-1)^jA_j(-X)$ and write
$L_m=\sum_{j=0}^{m-1}(-1)^jB_j(X)B_{m-1-j}(-X)$;
then the argument is an identity inside $\mathscr R$ and does not
require choosing an antiderivative $F$.

The normal-form projection respects factor degree, and replacing every
$X_j$ by $-X_j$ multiplies a degree-$\ell$ monomial by $(-1)^\ell$.
Applying the unique projection to the concomitant identity yields
$N_m(-X)=(-1)^mN_m(X)$, proving the parity assertion and
\eqref{dq:paritycount}. Finally substitute $A_m=N_m+DT_m$ into
\eqref{dq:Arecurrence} and use
$X_1DT_m=D(X_1T_m)-X_2T_m$. The result is
\[
 A_{m+1}=X_1N_m-X_2T_m+D(X_1T_m-N_m-DT_m),
\]
which proves \eqref{dq:pairrecurrence}.
\end{proof}

The values of $r_m$ for $m=2,\ldots,12$ are
\[
 1,1,2,2,3,3,5,6,8,9,13.
\]
These are formal upper bounds for the number of coordinates in
the particular Gamma normal form. No nonvanishing assertion for all
allowed coefficients is required. Some low-degree remainders are
\begin{align}
 N_4&=X_1^4-X_2^2,\notag\\
 N_5&=X_1^5-5X_1X_2^2,\notag\\
 N_6&=X_1^6-15X_1^2X_2^2+X_3^2,\label{dq:N4567}\\
 N_7&=X_1^7-35X_1^3X_2^2+7X_1X_3^2,\notag\\
 N_8&=X_1^8-70X_1^4X_2^2+28X_1^2X_3^2-7X_2^4-X_4^2.
 \notag
\end{align}
In particular the weight-six basis element $X_2^3$ is excluded by
parity. Its absence is a structural consequence, not a fitted numerical
cancellation.

\subsection{Appell transport of the reduced differential polynomials}

Let
\[
 P_j(Y)=j![t^j]\Gamma(1-t)e^{tY},\qquad
 C_m(x)=m![t^m]\frac{\Gamma(1-t)\Gamma(x)}{\Gamma(x-t)}.
\]
These are the established scalar and differential Gamma normalizers
from the incoming fifth-moment report. The operator $\partial_{X_1}$
commutes with $D$, and the repeated-largest normal-form space is
stable under $\partial_{X_1}$. Uniqueness gives
\[
 \partial_{X_1}N_m=mN_{m-1}.
\]
Consequently the canonical remainder of $C_m$ is obtained by expanding
$N_m$ as a polynomial in $X_1$ and replacing each power $X_1^j$ by
$P_j(\psi)$. Explicitly, with $K_j=N_j|_{X_1=0}$,
\begin{equation}\label{dq:Appelltransport}
 \mathcal N(C_m)=\sum_{j=0}^m\binom mj P_j(\psi)K_{m-j}.
\end{equation}
For example,
\begin{align}
 \mathcal N(C_6)&=P_6(\psi)-15P_2(\psi)(\psi')^2+(\psi'')^2,
 \label{dq:C6normal}\\
 \mathcal N(C_7)&=P_7(\psi)-35P_3(\psi)(\psi')^2
                                  +7P_1(\psi)(\psi'')^2.
 \label{dq:C7normal}
\end{align}
The primitive can be obtained by the same finite Gamma coefficient
convolution applied to the $T_j$; constant ambiguities do not affect
its endpoint functional.

Moreover, if $C_m=\mathcal N(C_m)+D\widetilde T_m$, then
\begin{equation}\label{dq:Appellendpointalgebra}
 \mathcal E(\widetilde T_m)\in
            \mathbb Q[\zeta(2),\ldots,\zeta(m-1)].
\end{equation}
It is homogeneous of weight $m-1$. To see this, put $y=\psi+\gamma$
and use
\[
 \Gamma(1-t)e^{-\gamma t}
  =\exp\left(\sum_{j\ge2}\frac{\zeta(j)}j t^j\right).
\]
The primitive is a differential polynomial in $y$ with these zeta
coefficients. Both $y(1)=0$ and the constant Laurent coefficient of
$y$ at zero are zero, so neither endpoint calculation introduces
$\gamma$. The finite formula \eqref{dq:finiteendpoint} proves
the assertion at every degree.

\subsection{The inherited spectral bridge in the new normal form}

Use the conventions
\[
 \frac1{\Gamma(1+u)}=\sum_{j\ge0}c_ju^j,\qquad
 \lambda_d=[u]\zeta(1+u,\underbrace{1,\ldots,1}_{d-1}),
 \quad d\ge1.
\]
The coefficient $[u]$ is a Laurent coefficient, not a derivative at
a pole. In particular $\lambda_1=-\gamma_1$ under the convention
$\zeta(1+u)=u^{-1}+\gamma-\gamma_1u+\cdots$.
The established Gamma integral bridge is
\begin{equation}\label{dq:inheritedCbridge}
 \int_0^1C_m(x)\,dx=m!(c_m-\lambda_{m-1}),\qquad m\ge2.
\end{equation}
Its complete integrand is continuous at zero. Combining this formula
with \eqref{dq:Appelltransport} and the exact endpoint functional
gives an all-degree finite recipe which isolates the single highest
height-one coefficient and specifies every remaining mixed integral.

For comparison, the equivalent inverse-Gamma bridge is
\begin{equation}\label{dq:inheritedAbridge}
 \mathcal L_m:=\operatorname{FP}\int_0^1 A_m(x)\,dx
 =m![t^m]\frac{\sin(\pi t)}{\pi t}
  -m!\sum_{d=1}^{m-1}(-1)^{m-d-1}c_{m-d-1}\lambda_d.
\end{equation}
This follows directly by multiplying the generating family $C_m$
by $1/\Gamma(1-t)$, retaining
$\operatorname{FP}\int C_0=1$ and
$\operatorname{FP}\int C_1=\gamma$, and using
$1/(\Gamma(1-t)\Gamma(1+t))=\sin(\pi t)/(\pi t)$.
Thus the new normal form changes the finite coordinate list and
retains the already established spectral normalization.

\section{A direct formula for coincident polygamma products}
\label{dq:bilinearsection}

The manuscript already gives general translated polygamma products
in first Hurwitz derivatives, in particular its formula
\texttt{shj:polyclosurezeta}. The following coincident, original-coordinate
finite part has the corresponding ordinary-zeta form. A direct proof
is supplied because a collision limit cannot in general be moved through
a finite-part prescription without checking its local normalization.

Put $H_0=0$ and $H_d=\sum_{j=1}^d1/j$.
For nonnegative integers $r,s$ with $d=r+s\ge1$, define
$\sigma_r=(-1)^{r+1}$ and $\sigma_s=(-1)^{s+1}$.

\begin{theorem}[Coincident bilinear polygamma finite parts]
\label{dq:bilinear}
In the original coordinate $x$,
\begin{align}
 &\operatorname{FP}\int_0^1\psi^{(r)}(x)\psi^{(s)}(x)\,dx\notag\\
 &\quad=d!\left\{(\sigma_r+\sigma_s)\zeta'(d+1)
 +\left[\sigma_r(H_d-H_s)+\sigma_s(H_d-H_r)\right]
                                      \zeta(d+1)\right\}.
 \label{dq:bilinearformula}
\end{align}
In particular, for every $r\ge1$,
\begin{equation}\label{dq:squareformula}
 \operatorname{FP}\int_0^1(\psi^{(r)}(x))^2\,dx
 =2(-1)^{r+1}(2r)!
    \left\{\zeta'(2r+1)+(H_{2r}-H_r)\zeta(2r+1)\right\}.
\end{equation}
\end{theorem}

\begin{proof}
At the pole $x=-n$, write $v=x+n$ and
\[
 \psi(-n+v)=-v^{-1}+a_0(n)+\sum_{k\ge1}a_k(n)v^k,
\]
where $a_0(n)=H_n-\gamma$ and
$a_k(n)=H_n^{(k+1)}+(-1)^{k+1}\zeta(k+1)$ for $k\ge1$.
In every case $\Delta a_k(n)=n^{-k-1}$.
Let $p_j(n)$ denote the coefficient of $v^{-j}$ in the product.
There is a constant highest coefficient at order $d+2$, and for
$1\le j\le d+1$ the other coefficients are
\begin{equation}\label{dq:bilinearpoles}
 p_j(n)=(d+1-j)!a_{d+1-j}(n)
 \left\{\mathbf1_{j\le r+1}\frac{\sigma_r r!}{(r+1-j)!}
       +\mathbf1_{j\le s+1}\frac{\sigma_s s!}{(s+1-j)!}\right\}.
\end{equation}
This is just multiplication of each leading pole by the other's
analytic Taylor series. In particular
\[
 \Delta p_1(n)=d!(\sigma_r+\sigma_s)n^{-d-1}.
\]

Here are the convergence and normalization details. The normalized
Mittag--Leffler expansion is
\[
 f(z)=P_0(z)+\kappa+
           \sum_{n\ge1}\{P_n(z)-P_n(0)\},\qquad
 P_n(z)=\sum_j\frac{p_j(n)}{(z+n)^j},
\]
where $\kappa$ is the local Laurent constant at zero. Its normal
convergence follows from $p_j(n)=O(1+\log n)$ and one subtraction.
The whole-circle contour argument eliminates the entire remainder:
integrate $z(f(w)-P_0(w))/(w(w-z))$ on $|w|=N+1/2$.
Digamma reflection and the differentiated cotangent reflection formulas
bound all factors by $O(1+\log N)$ on this circle, so the integral
tends to zero. Moreover
\[
 \kappa-\sum_{n=1}^NP_n(0)
       =\frac1{2\pi}\int_{-\pi}^{\pi}f((N+1/2)e^{i\theta})\,d\theta
       \longrightarrow0.
\]
For the last limit, outside an arc of length $O(\log N/N)$ at the
negative real axis, at least one differentiated digamma factor tends
to zero by the sectorial asymptotic and reflection. On the exceptional
arc the uniform logarithmic bound times its length tends to zero.
This uses $d\ge1$; the undifferentiated square is a separate case.

Integrate the normal expansion to one before passing to the cutoff.
Finite summation by parts in its higher-pole terms gives
\[
 \sum_{j\ge2}\sum_{n\ge1}
              \frac{\Delta p_j(n)}{(j-1)n^{j-1}},
\]
and the simple poles give $-\sum_{n\ge1}\Delta p_1(n)\log n$.
All terminal terms vanish: $p_j(N)/(N+1)^{j-1}\to0$ for $j\ge2$;
also $p_1(N)\log(N+1)\to0$, since it is zero when $d$ is odd,
and proportional to the decaying tail $a_d(N)$ when $d$ is even.
The grouped sums are absolutely convergent. They are precisely the
original-coordinate finite part, by integration before constant-term
extraction at zero.

The logarithmic sum is $d!(\sigma_r+\sigma_s)\zeta'(d+1)$.
For the higher poles use the finite identity
\begin{equation}\label{dq:binomialharmonic}
 r!\sum_{h=1}^r\frac{(d-h)!}{(r-h)!h}
                 =d!(H_d-H_{d-r}),\qquad0\le r\le d.
\end{equation}
After division by $d!$, this identity is obtained by integrating
\[
 \sum_{h=1}^r\binom rh t^{h-1}(1-t)^{d-h}
       =\frac{(1-t)^{d-r}-(1-t)^d}{t}
\]
from zero to one. The beta integrals on the left and the finite
geometric sum on the right give exactly \eqref{dq:binomialharmonic}.
Applying it once to each term in \eqref{dq:bilinearpoles} gives
the zeta coefficient in \eqref{dq:bilinearformula}. Setting $s=r$
proves \eqref{dq:squareformula}.
\end{proof}

The cases needed below are
\begin{equation}\label{dq:twosquares}
 \operatorname{FP}\int_0^1(\psi')^2\,dx=2\zeta(3)+4\zeta'(3),
 \qquad
 \operatorname{FP}\int_0^1(\psi'')^2\,dx=-28\zeta(5)-48\zeta'(5).
\end{equation}

\section{The canonical sixth and seventh digamma identities}
\label{dq:sixthsection}

Write $z_j=\zeta(j)$ in this section and put $y=\psi+\gamma$.
The sixth scalar polynomial is explicitly
\begin{align}
 P_6(\psi)={}&y^6+15z_2y^4+40z_3y^3
          +(45z_2^2+90z_4)y^2\notag\\
 &+(120z_2z_3+144z_5)y
     +15z_2^3+90z_2z_4+40z_3^2+120z_6.
 \label{dq:P6}
\end{align}
Define the one retained mixed moment
\begin{equation}\label{dq:M6}
 \widetilde M_6=\operatorname{FP}\int_0^1
                          (\psi(x)+\gamma)^2(\psi'(x))^2\,dx.
\end{equation}
This is a fully specified original-coordinate finite part, with an
ordinary convergent subtracted representation obtained by its finite
Laurent principal part.
More explicitly,
\begin{equation}\label{dq:M6principal}
 y^2(\psi')^2=x^{-6}-2z_3x^{-3}-z_4x^{-2}
                 +6(z_2z_3-z_5)x^{-1}+O(1).
\end{equation}
Consequently
\begin{align}
 \widetilde M_6={}&-\frac15+z_3+z_4\notag\\
 &+\int_0^1\left\{y^2(\psi')^2-x^{-6}+2z_3x^{-3}
         +z_4x^{-2}-6(z_2z_3-z_5)x^{-1}\right\}\,dx.
 \label{dq:M6ordinary}
\end{align}
The complete integrand in braces has a removable singularity at zero.

\begin{theorem}[A canonical sixth-degree bridge with one mixed moment]
\label{dq:canonicalsixth}
With the conventions already fixed,
\begin{align}
 \boxed{\begin{aligned}
 \operatorname{FP}\int_0^1P_6(\psi(x))\,dx
 ={}&720(c_6-\lambda_5)+15\widetilde M_6\\
 &+60z_2\zeta'(3)+48\zeta'(5)+90z_2z_3+270z_5.
 \end{aligned}}
 \label{dq:canonicalsixthformula}
\end{align}
Thus this scalar Gamma--Appell combination has only the highest
height-one coefficient $\lambda_5$ and one specified mixed moment,
apart from ordinary constants and ordinary zeta derivatives. This is
not an elimination or an independence theorem for $\widetilde M_6$.
\end{theorem}

\begin{proof}
For the inverse-Gamma polynomial, the exact primitive from the normal
algorithm is
\begin{align}
 T_6={}&-3\psi^5+20\psi^3\psi'-15\psi^2\psi''
       -15\psi(\psi')^2+6\psi\psi'''+9\psi'\psi''-\psi''''.
 \label{dq:T6}
\end{align}
Direct differentiation, or the finite recurrence, gives
\[
 A_6=\psi^6-15\psi^2(\psi')^2+(\psi'')^2+T_6'.
\]
For $C_6$, form the polynomial
\begin{align}
 \widetilde T_6={}&T_6[y]
  +15z_2(-2y^3+4y\psi'-\psi'')\notag\\
  &+40z_3(-\tfrac32y^2+\psi')-(45z_2^2+90z_4)y,
 \label{dq:TC6}
\end{align}
where $T_6[y]$ means \eqref{dq:T6} with every undifferentiated
$\psi$ replaced by $y$ and its derivatives unchanged. Then
\begin{equation}\label{dq:C6exact}
 C_6=P_6(\psi)-15(y^2+z_2)(\psi')^2+(\psi'')^2
                              +\widetilde T_6'.
\end{equation}
The two endpoint values are
\begin{align}
 \widetilde T_6(1)&=52z_2z_3+24z_5,\notag\\
 \operatorname{CT}_{x=0}\widetilde T_6(x)
                  &=112z_2z_3+266z_5,\notag\\
 \mathcal E(\widetilde T_6)&=-60z_2z_3-242z_5.
 \label{dq:C6endpoints}
\end{align}
They follow by substituting
$y=-x^{-1}+z_2x-z_3x^2+z_4x^3-z_5x^4+O(x^5)$,
or directly from \eqref{dq:finiteendpoint}. Apply finite parts to
\eqref{dq:C6exact}, use \eqref{dq:inheritedCbridge}, and then use
both evaluations in \eqref{dq:twosquares}. Every term in
\eqref{dq:canonicalsixthformula} follows, including the endpoint
contribution and the $z_2\zeta'(3)$ correction.
\end{proof}

For a fully expanded ordinary coefficient,
\begin{align}
 720c_6={}&\gamma^6-15\gamma^4z_2+40\gamma^3z_3
       +45\gamma^2z_2^2-90\gamma^2z_4\notag\\
 &-120\gamma z_2z_3+144\gamma z_5
       -15z_2^3+90z_2z_4+40z_3^2-120z_6.
 \label{dq:c6}
\end{align}

\subsection{The raw sixth power}

Let $Q_6=\operatorname{FP}\int_0^1\psi^6$ and
$M_6=\operatorname{FP}\int_0^1\psi^2(\psi')^2$.
The endpoint difference of \eqref{dq:T6} is
\begin{equation}\label{dq:rawsixthendpoint}
 \mathcal E(T_6)=60\gamma^3z_2-150\gamma^2z_3-105\gamma z_4
                              +210z_2z_3-242z_5.
\end{equation}
Combining this with \eqref{dq:inheritedAbridge} gives the complete
raw-power formula
\begin{align}
 Q_6={}&15M_6-720\lambda_5+720\gamma\lambda_4
          -360(\gamma^2-z_2)\lambda_3\notag\\
 &+120(\gamma^3-3\gamma z_2+2z_3)\lambda_2\notag\\
 &-30(\gamma^4-6\gamma^2z_2+8\gamma z_3
                       +3z_2^2-6z_4)\lambda_1\notag\\
 &-135z_6+48\zeta'(5)-60\gamma^3z_2+150\gamma^2z_3
             +105\gamma z_4-210z_2z_3+270z_5.
 \label{dq:rawsixth}
\end{align}
The term $-135z_6=-\pi^6/7$ is the entire contour constant of
the sixth inverse-Gamma bridge. It is not an endpoint correction.

\subsection{The next odd degree}

The algorithm also gives a compact seventh-degree extension. Define
\[
 W_7=\operatorname{FP}\int_0^1P_3(\psi)(\psi')^2\,dx,
 \qquad
 K_7=\operatorname{FP}\int_0^1P_1(\psi)(\psi'')^2\,dx.
\]
Then
\begin{equation}\label{dq:canonicalseventh}
 \boxed{\displaystyle
 \operatorname{FP}\int_0^1P_7(\psi)\,dx
 =5040(c_7-\lambda_6)+35W_7-7K_7+231z_3^2+2380z_6.}
\end{equation}
Here $P_1(\psi)=y$ and $P_3(\psi)=y^3+3z_2y+2z_3$.
For reproducibility, a primitive for $A_7-N_7$ is
\begin{align}
 T_7={}&-\tfrac72\psi^6+35\psi^4\psi'-35\psi^3\psi''
        -\tfrac{105}2\psi^2(\psi')^2+21\psi^2\psi'''\notag\\
 &+63\psi\psi'\psi''-7\psi\psi''''+14(\psi')^3
        -14\psi'\psi'''-\tfrac{21}2(\psi'')^2+\psi'''''.
 \label{dq:T7}
\end{align}
In addition to $T_6$ already displayed, the lower primitives needed
for the Gamma convolution are
\begin{align*}
 T_5&=-\tfrac52\psi^4+10\psi^2\psi'-5\psi\psi''
                         -\tfrac52(\psi')^2+\psi''',\\
 T_4&=-2\psi^3+4\psi\psi'-\psi'',\\
 T_3&=-\tfrac32\psi^2+\psi',\qquad T_2=-\psi.
\end{align*}
The primitive for $C_7-\mathcal N(C_7)$ is
\begin{align*}
 \widetilde T_7={}&T_7[y]+21z_2T_5[y]+70z_3T_4[y]\\
  &+(105z_2^2+210z_4)T_3[y]+(420z_2z_3+504z_5)T_2[y].
\end{align*}
The finite endpoint algorithm gives
\[
 \mathcal E(\widetilde T_7)
 =-448z_2z_4-231z_3^2-1596z_6
 =-231z_3^2-2380z_6.
\]
Combining it with \eqref{dq:C7normal} and the inherited Gamma
bridge proves \eqref{dq:canonicalseventh}. The certificate contains
the full polynomial identity, so this seventh-degree formula does
not rely on a numerical extrapolation from degree six.

\subsection{Proof status, verification, and remaining question}

The accompanying exact program checks every partition at every weight
through twelve: it verifies the decomposition identity, the vanishing
normal form of every total derivative, the bidegree count, the adjoint
concomitant identity, and the parity restriction. It also independently
expands the local Laurent polynomials for the displayed sixth and
seventh endpoint constants. These finite checks corroborate the
proofs; the terminating algorithm and the dimension theorem themselves
hold at all degrees by the arguments above.

This resolves the formal integration-by-parts classification requested
in the incoming programme and gives complete degree-six and degree-seven
applications. The next separate task is to determine the image of this
formal quotient in an explicitly specified family of harmonic spectral
coordinates. Its kernel may contain additional special-function
relations. Neither the partition count nor the parity count can be
used to assert arithmetic independence of the remaining constants.

% Suggested bibliography entry for the surrounding article:
% \bibitem{OlverShakiban1978}
% P. J. Olver and C. Shakiban, A resolution of the Euler operator. I,
% Proc. Amer. Math. Soc. 69 (1978), 223--229.
% DOI: 10.1090/S0002-9939-1978-0486822-8.
% Primary author-hosted PDF: https://www-users.cse.umn.edu/~olver/s_/reo.pdf

% END sections/06_differential.tex


% BEGIN sections/07_audit_research.tex
\section{Source audit, proposed status changes, and further research}
\label{sec:audit-research}

\subsection{What should change in the existing programme}

The audit was directed at the statements used or continued in this
article. It was not a line-by-line certification of every chapter in the
large collective manuscript. No mathematical counterexample or invalid
proof step was established in those source statements. The principal
proposed changes are strengthened theorems and more precise boundaries
between formal and arithmetic conclusions.

\begin{enumerate}[leftmargin=*,label=\arabic*.]
\item \textbf{Split the optimal-prefix conjecture by scope.}
In \cite[\texttt{conj:firstq}]{CanonicalFifth}, replace the $p=2$,
$p=4$, and $p=6$ parts by Theorems~\ref{thm:rank-two}, \ref{rf:main},
and \ref{rs:main},
including their determinants and Corollary~\ref{cor:ordinary-two}.
Retain the diagonal theorem already proved there. Theorem
\ref{thm:general-small-determinant} gives an explicit fixed-size
reduction for all even pairs. Leave unequal even pairs $8\le p<q$
as the remaining conjecture. Finite checks at those orders should continue to be identified
as checks, with no inferred all-order conclusion.

\item \textbf{Extend the directional-zeta chapter through cubic order.}
Append Theorem~\ref{tc:main} and the two explicit cancellations
\eqref{tc:two-ray-example} and \eqref{tc:example} after the inherited
quadratic theory. Preserve the precise admissible-ray condition and
the displayed polar subtraction. A global third Taylor tensor of $T$
at the origin is not defined; the theorem concerns specified one-variable
restrictions. The kernel theorem should be attached to the exact local
restrictions it classifies.

\item \textbf{Replace the unstructured higher-power task by a proved
algorithm and a smaller arithmetic task.}
Theorems~\ref{dq:normalform} and \ref{dq:parity} settle the formal
integration-by-parts classification in the chosen differential-polynomial
space. Equations~\eqref{dq:canonicalsixthformula} and
\eqref{dq:canonicalseventh} give explicit applications. What remains is
the spectral evaluation of their specified mixed moments, not the
existence of an integration-by-parts procedure.

\item \textbf{Preserve the two finite-part normalizations.}
Centered tail moments use continuation of $D_{p,q}(s)$, while singular
integrals use the constant term in the original coordinate $x$.
The proofs state why the particular ordinary subtractions recover those
values. An unannounced translation of the cutoff or omission of
$\mathcal E(T)$ could alter constants. The special binomial omission
rule in \eqref{eq:entries-foundation} must also remain explicit: restoring
a forbidden lower-index-zero term changes the matrix under discussion.

\item \textbf{Keep the priority and independence statements narrow.}
The Bernoulli Hankel determinant and differential-polynomial exactness
machinery have classical antecedents. The height-one coefficient bridge
is inherited. The new content lies in the mixed determinant, the
resulting exact cancellations, the cubic directional reduction, and the
explicitly normalized higher-power applications. Formal coordinate counts
do not prove independence of zeta values or of the remaining mixed
integrals. No historical priority assertion beyond the inspected material
is needed for these mathematical contributions.
\end{enumerate}

The manuscript's short $S_6$ identity and its revised $S_8$ candidate
remain conjectural. The older rejected $S_8$ vector is a different
statement; it must not be substituted for the current candidate. The
present results neither prove nor refute those period identities
\cite[README and editorial ledger]{Repo}.

\subsection{A proposed programme of exact research questions}

The following questions are selected for identities, reduction theorems,
and explicit primitives. They are ordered approximately by their
connection to the results above, rather than by an unproved prediction
of difficulty.

\begin{enumerate}[leftmargin=*,label=\textbf{R\arabic*.}]
\item \textbf{Prove the remaining unequal-order prefix conjecture.}
For $p=2a<2b=q$ with $a\ge4$, use the explicit $a$-square matrix
in Theorem~\ref{thm:general-small-determinant} to finish the remaining
nonvanishing problem. The shifted-Hankel part is already known to be
nonsingular, and the weighted Bernoulli reduction has degree at most
$2a-1$. The next task is to evaluate the residual determinant and find
an arithmetic obstruction to its vanishing, as in the $a=1,2,3$ proofs. The exact $(a,b)=(4,9)$ example
in Section~\ref{rs:section} rules out the naive uniform valuation $-a$;
any more general argument must accommodate that obstruction.
A proof must control
nonvanishing of the final rational factor, not just display a product of
unknown minors.

\item \textbf{Determine every reconstruction row in closed form.}
Theorem~\ref{thm:rational-row-two} gives the rational coordinate $Z_0$.
An explicit formula for every row of $M_b^{-1}$ would express each even
zeta coordinate through the same optimal moment prefix. Orthogonal
polynomial kernels followed by the rank-one inverse formula offer a
concrete route. Seek coefficient formulas with finite binomial sums and
prove their normalization directly. The result would produce systematic
ordinary harmonic-tail identities with prescribed $\zeta(2k)$ values.

\item \textbf{Extend the rational tail cancellations to shifted or
character-weighted tails.}
Replace the ordinary tails by Hurwitz tails with a rational shift, or
by tails weighted by Dirichlet characters. Determine which centering and
parity permit an exact polynomial subtraction, and which generalized
Bernoulli determinants replace $M_b$. The half-integer vanishing
$\zeta(-2r,1/2)=0$ is a special ingredient of the present proof; in a new
shift it must be replaced by an explicitly computed finite correction.

\item \textbf{Find finite polylogarithmic evaluations associated with
the boundary family.}
Equation~\eqref{eq:polylog-boundary} is an exact endpoint identity for
an explicit finite differential operator. Determine whether its numerator
admits useful integrated forms at interior algebraic arguments, with
depth reduction proved by differential identities and base values.
The test is an equality of functions or a rigorously anchored primitive,
not a numerical integer relation at a single argument.

\item \textbf{Determine the arithmetic meaning of the second cubic
Tornheim coordinate.}
The constant $\kappa$ is fixed by $J_{AA}(0)$ in the canonical
subtraction, and \eqref{tc:kappa-Euler} already expresses it through
the strict Euler double-zeta ray $\zeta(2t,t)$. Starting from that
proved representation, compare it with the inherited cubic digamma
invariant and seek a reduction. The concrete question is whether
$\kappa$ belongs to a specified algebra generated by ordinary zeta
derivatives and $\Omega$, or whether a further higher-depth coordinate
is required. A coordinate change alone does not answer that question.

\item \textbf{Compute the quartic and quintic ray cancellation laws.}
Theorem~\ref{tc:kernel} gives formal ambiguity dimensions $3$ and $6$
at degrees four and five. Construct explicit homogeneous direction
functions for these spaces, derive the ordinary-zeta inhomogeneous part,
and find short integer-weighted combinations of positive rays that
annihilate all undetermined coefficients. These are finite symbolic
problems once the analytic polar normalization has been fixed, but the
coefficient calculation must be independently checked.

\item \textbf{Add further Tornheim functional relations to the local
kernel and compute the exact reduction.}
The kernel theorem deliberately uses a specified subset of the known
relations. Linearize additional functional equations from the classical
Tornheim literature at the same origin, keeping their Gamma and trigonometric
factors to the required order. Determine whether they reduce the cubic
two-coordinate space or only change its basis. At higher degrees, compute
the quotient of the kernel by these additional restrictions before making
claims about necessary arithmetic data.

\item \textbf{Evaluate the retained sixth mixed moment in an explicitly
stated spectral family.}
The quantity $\widetilde M_6$ has a convergent subtracted integral in
\eqref{dq:M6ordinary}. Use its pole principal parts, rather than numerical
fitting, to express it through ordered harmonic zeta derivatives. Then
test whether the resulting combination reduces within a fixed family of
height-one, depth-two, or Tornheim coordinates. The desired theorem is an
explicit evaluation or a proved reduction in that family; the present
formula makes no assertion that $\widetilde M_6$ is irreducible.

\item \textbf{Determine the exact image of the differential quotient.}
The partition theorem classifies formal integration by parts. Define the
map from this quotient, together with the endpoint functional, to the
harmonic spectral coordinates supplied by the pole expansion. Compute
its kernel at each weight and identify structural generators for the
additional relations. The square formula is an example of a relation
available after substitution into the digamma function that is not an
identity in the free differential-polynomial quotient.

\item \textbf{Combine higher Stieltjes coefficients with the normal
form and parameter primitives.}
The fixed Mellin family in Proposition~\ref{prop:height-one-foundation}
gives every $\lambda_{d,j}$ and ordinary $\gamma_j$ at depth one.
Introduce the appropriate logarithmic parameter derivatives into the
Gamma family, determine their differential normal forms, and retain the
complete endpoint terms. A useful target is a closed recurrence for
primitive families whose spectral derivatives isolate one prescribed
$\lambda_{d,j}$, followed by analogous formulas for Hurwitz
$\gamma_j(a)$ with a fixed shift convention.

\item \textbf{Explain the coefficient patterns of the Gamma normal
forms.}
The Appell property already proves, for example, that the coefficient of
$X_1^{m-4}X_2^2$ is $-\binom m4$: its undifferentiated core is the
$-X_2^2$ term in $N_4$. The remaining coefficient problem is therefore
concentrated in $K_j=N_j|_{X_1=0}$. Prove general formulas for families
of repeated-largest partitions in these cores using the finite recurrence.
Distinguish allowed parity coordinates from coordinates whose coefficient
actually vanishes; an all-degree support theorem would sharpen the
current formal upper count.

\item \textbf{Return to the short Gaussian identities with a fixed
period target.}
For $S_6$ and revised $S_8$, retain the repository's exact frozen
vectors and the existing shuffle equivalences. Seek a new functional or
iterated-integral identity whose boundary evaluation is exactly that
vector. A nonzero projection in a chosen formal relation space does not
refute the numerical period identity, and a small computed residual does
not prove it. The new determinant and differential methods may suggest
useful auxiliary identities, but no direct implication is established here.
\end{enumerate}

Each question specifies the missing mathematical step. In particular,
the status of a constant evaluation should change only when an analytic
identity or an exact algebraic reduction with justified boundary values
has been supplied.

% END sections/07_audit_research.tex


% BEGIN sections/08_verification.tex
\section{Reproducibility and the limits of computation}
\label{sec:verification}

The accompanying package contains the complete typeset article, a single
flattened \TeX\ source, the modular sources, eleven verification programs,
their recorded JSON results, and an integration guide. The proofs in the
article establish the uniform statements. The programs provide exact
finite certificates and independent numerical diagnostics at explicitly
specified orders.

\subsection{Exact finite certificates}

The order-two verifier constructs the matrix directly from the inherited
binomial formula, without defining it by the proposed factorization.
For every $b=2,\ldots,14$ it checks the determinant, the rank-one
factorization, the Pad\'e inverse-row contraction, the rational extraction
row, the shifted Hankel product, and the final factorial cancellation.
It also checks the Pad\'e coefficient identities through the corresponding
orders and reconstructs the ordinary-series examples at $b=2,3,4$.

The order-four verifier checks the symbolic perturbation contractions,
the resulting determinant polynomial, and its exact parity calculation.
It also constructs the original source matrices for $b=3,\ldots,14$,
including the exceptional smallest pair $(4,6)$. The general shifted-Hankel
verifier checks all 45 pairs $1\le a\le b\le9$, comparing the source
entries with the factorization, the determinant signs, the size-$a$
reduction, and the weighted Bernoulli sums. The exact ranges supplement
the uniform proofs of both new sections.

The order-six program constructs the normalized $3\times3$ matrix with
an indeterminate $b$. It proves the determinant normalization, constructs
its numerator in the integer polynomial ring, and verifies the modulo-four
specialization used in the nonvanishing proof. Direct matrices are also
checked at $b=4,\ldots,12$, and the order-eight counterexample to the
uniform valuation pattern is computed exactly. A separate short program
checks the explicit $(4,6)$ ordinary tail identity. These computations
are finite polynomial certificates, not interpolation from integer samples.

The Tornheim symbolic verifier multiplies the Gamma and zeta Taylor
series in exact rational arithmetic. It verifies both elementary
restrictions on $J$, the full cubic direction formula, and both printed
cancellation identities. A separate construction builds the linear
restriction matrices for homogeneous polynomial germs through degree ten
and checks the stated kernel dimensions. An additional independent script
derives the three-ray cancellation directly from raw series and the two
elementary constraints; it also detects a deliberately changed coefficient.

The differential verifier checks every partition monomial at every
weight through twelve. It verifies its exact decomposition into normal
form and derivative, the zero normal form of every derivative column,
the factor-degree counts, the adjoint concomitant, and the Gamma parity.
It constructs the full degree-six and degree-seven polynomial identities,
including the canonical endpoint constants. The Laurent principal part
of the retained sixth mixed integral is checked symbolically as well.

The checks over finite ranges confirm the formulas within those
ranges. They do not by themselves prove the uniform statements.
The identities certified with indeterminate $b$, including the
integer-polynomial certificate in Section~\ref{rs:section}, are exact
finite algebraic steps in the proofs. The uniform conclusions also
use the Pad\'e coefficient argument, the Bernoulli congruences,
holomorphic divisibility, and the terminating partition argument,
as applicable. No external proof assistant was used to formalize
the article's analytic proofs.

\subsection{Independent numerical diagnostics}

Table~\ref{tab:numerics} reports the main recorded comparisons. Every
entry is floating-point evidence, not an interval enclosure. The numerical
evaluators do not assume the identities they are checking.

\begin{table}[htbp]
\centering\small
\caption{Recorded numerical comparisons; residuals are diagnostic.}
\label{tab:numerics}
\begin{tabular}{@{}p{.33\textwidth}p{.36\textwidth}>{\raggedleft\arraybackslash}p{.24\textwidth}@{}}
\toprule
Identity & Independent evaluation & Absolute residual\\
\midrule
Ordinary $(2,4)$ tail sum & 100-digit tails; 100 initial terms and 24 asymptotic terms
 & $9.93\cdot10^{-70}$\\[5pt]
Ordinary $(2,6)$ tail sum & Same method, independently formed weight polynomial
 & $3.20\cdot10^{-65}$\\[5pt]
Ordinary $(2,8)$ tail sum & 100 digits; 120 initial terms and 32 asymptotic terms
 & $3.28\cdot10^{-72}$\\[5pt]
Three cubic Tornheim rays & Subtracted Mellin evaluation and odd-part extrapolation
 & $6.30\cdot10^{-33}$\\[5pt]
Canonical sixth digamma formula & Direct Laurent subtraction and separate height-one Mellin integral
 & $5.63\cdot10^{-51}$\\[5pt]
$(\psi')^2$ and $(\psi'')^2$ & Direct Laurent subtraction and ordinary quadrature
 & $\begin{gathered}1.23\cdot10^{-90},\\1.77\cdot10^{-88}\end{gathered}$\\
\bottomrule
\end{tabular}
\end{table}

For the harmonic sums, the exact polynomial subtraction removes every
nonnegative asymptotic power; the remaining inverse powers are summed
by ordinary Hurwitz zeta values to accelerate a convergent series.
For the independent $(2,8)$ check, increasing the precision and the
tail parameters changed the residual from approximately
$4.20\cdot10^{-52}$ to $3.28\cdot10^{-72}$.

The Tornheim evaluator integrates the Jonqui\`ere expansion below $x=1$
and sums an incomplete-Gamma tail above $x=1$. With 50-digit arithmetic,
series index 48, tail index 110, and seven odd-part steps starting at
$h=0.0008$, it gives
\[
 W_{1,2,3}'''(0)+4W_{2,3,1}'''(0)-5W_{3,1,2}'''(0)
 \approx-451.3535195684862560644674145675405.
\]
The code does not use the proposed directional formula, fit $\Omega$
or $\kappa$, or derive its values from the digamma moment.

For the sixth identity, direct finite-part quadrature gives
\begin{align*}
 \FP\int_0^1P_6(\psi(x))\,dx
  &\approx20.8226131542578224267340551051333,\\
 \widetilde M_6
  &\approx-3.40244508656918661747498017448045.
\end{align*}
The separate Mellin integral uses an explicit lower cutoff $10^{-60}$;
the omitted piece is the limiting numerical effect in the displayed
comparison. Its size was not promoted to a certified error bound. The
formula itself follows from exact polynomial identities, exact endpoints,
and the proved ordinary Gamma bridge.

\subsection{Build and integration procedure}

Run \src{python build.py} to regenerate the flattened source and PDF.
The build uses \src{pdflatex} and resolves cross-references in repeated
passes. Run \src{python verification/verify_all.py} for the exact
certificates. The optional flag \src{--numerical} also runs the numerical
diagnostics; the Mellin computation is the longest job.
Required Python packages and the versions used are recorded in the
package. The JSON files contain parameters and full recorded output,
so a future replay can distinguish a change in arithmetic, a change in
numerical settings, and a change in the mathematical claim.

The integration guide maps each contribution to its source question and
names the labels that may be imported into the manuscript. It proposes
status edits without changing the pinned source checkout. The article is
intended to be read and checked as a research continuation before those
edits are adopted.

% END sections/08_verification.tex


% BEGIN references.tex
\begin{thebibliography}{99}

\bibitem{Repo}
V.~Reshetnikov and contributors,
\emph{ProveIt: Polylogarithms and their Arithmetic Bridges},
collective manuscript and incoming research, inspected at commit
\texttt{f7e35be357654198c1546fd40f363694c5bcc2cc} (11 October 2026).
Manuscript: \url{https://github.com/VladimirReshetnikov/ProveIt/tree/f7e35be357654198c1546fd40f363694c5bcc2cc/Analysis/Polylogarithms/docs/manuscript}.
Incoming: \url{https://github.com/VladimirReshetnikov/ProveIt/tree/f7e35be357654198c1546fd40f363694c5bcc2cc/docs/incoming}.

\bibitem{CanonicalFifth}
ProveIt incoming research,
\emph{Canonical Digamma and Fifth Moment}, 11 October 2026.
Archive \src{docs/incoming/ProveIt_Canonical_Digamma_Fifth_Moment_2026-10-11.zip}
at the pinned repository commit \cite{Repo}; especially
\src{sections/02_laurent.tex}, \src{03_gamma.tex},
\src{07_tail_rank.tex}, and \src{09_questions.tex}.
The archive is the mathematical baseline for the optimal-prefix
conjecture and the Gamma integral normalization used here.

\bibitem{HigherPowers}
ProveIt incoming research,
\emph{Higher Digamma Powers and Ordered Gamma Sums}, 11 October 2026.
Archive \src{docs/incoming/ProveIt_Higher_Digamma_Powers_2026-10-11.zip}
at the pinned repository commit \cite{Repo}; especially its sections
on differential polynomials and further exact research questions.

\bibitem{DirectionalPrior}
ProveIt research continuation,
\emph{Coincident and Directional Zeta}, as present at the pinned
repository commit \cite{Repo}.
Directory
\src{Analysis/Polylogarithms/docs/reports/stieltjes-correlation-calculus/continuations/coincident-directional-zeta};
especially \src{sections/02_cubic.tex} and \src{sections/06_rays.tex}.

\bibitem{DLMF}
\emph{NIST Digital Library of Mathematical Functions}.
Sections 4.36, 5.15, 24.10, 25.11, and 25.12.
\url{https://dlmf.nist.gov/}. Accessed 11 October 2026.

\bibitem{DilcherJiu2020}
K.~Dilcher and L.~Jiu,
\emph{Hankel determinants of sequences related to Bernoulli and
Euler polynomials}, author preprint, 2020.
\url{https://arxiv.org/abs/2007.09821}.
Section 7.3, equation (7.12), records the normalized shifted-Bernoulli
Hankel evaluation and attributes it to Krattenthaler.

\bibitem{OlverShakiban1978}
P.~J. Olver and C.~Shakiban,
\emph{A resolution of the Euler operator. I},
Proc. Amer. Math. Soc. \textbf{69} (1978), 223--229.
\url{https://doi.org/10.1090/S0002-9939-1978-0486822-8}.
Author-hosted text:
\url{https://www-users.cse.umn.edu/~olver/s_/reo.pdf}.

\bibitem{Young2024}
P.~T. Young,
\emph{Series of height one multiple zeta functions},
Integers \textbf{24} (2024), Paper A43, 14 pp.
\url{https://doi.org/10.5281/zenodo.11221606}.
\url{https://math.colgate.edu/~integers/y43/y43.pdf}.
Section 4 gives the height-one Laurent-coefficient context. The present
article fixes its own alternating Stieltjes sign convention explicitly.

\bibitem{Sesma}
J.~Sesma,
\emph{The Roman harmonic numbers revisited}, author preprint, 2017.
\url{https://arxiv.org/abs/1702.03718}.

\bibitem{TC:BD}
J.~M.~Borwein and K.~Dilcher,
\emph{Derivatives and fast evaluation of the Tornheim zeta function},
Ramanujan J. \textbf{45} (2018), 413--432.
\url{https://doi.org/10.1007/s11139-017-9890-9}.
Author preprint, titled \emph{Derivatives and fast evaluation of the
Witten zeta function}:
\url{https://carmamaths.org/resources/jon/omega3.pdf}.

\bibitem{TC:Romik}
D.~Romik,
\emph{On the number of $n$-dimensional representations of
$\operatorname{SU}(3)$, the Bernoulli numbers, and the Witten zeta
function}, Acta Arith. \textbf{180} (2017), 111--159.
\url{https://arxiv.org/abs/1503.03776}.

\bibitem{TC:Onodera2014}
K.~Onodera,
\emph{A functional relation for Tornheim's double zeta functions},
Acta Arith. \textbf{162} (2014), 337--354.
\url{https://doi.org/10.4064/aa162-4-2}.
Author preprint: \url{https://arxiv.org/abs/1211.1480}.

\bibitem{TC:Onodera2021}
K.~Onodera,
\emph{On multiple Hurwitz zeta function of Mordell--Tornheim type},
Int. J. Number Theory \textbf{17} (2021), 2327--2360.
\url{https://doi.org/10.1142/S1793042121500913}.
The comparison made here uses the published abstract and bibliographic
scope; the subscription full text was not available for a complete
formula-by-formula priority audit.

\end{thebibliography}

% END references.tex

\end{document}
